% Quantum Gravity Dynamics — new Chapter 4, independent research manuscript
% This is NOT an edit to the historical Chapter 4.
% Original 19 October 2026; QGD-097 updated 26 October 2026.
% QGD-098 surgical referee-verification insertion 26 October 2026.
% QGD-101 surgical integration of research notes QGD-099/100, 26 October 2026.
% QGD-122 insertion, 26 October 2026: spanning representations of the chapters' backgrounds.
% QGD-121 insertion, 26 October 2026: degree-of-freedom count of the QGD action.
% QGD-120 insertion, 26 October 2026: the complete QGD action.
% QGD-116 cross-reference, 26 October 2026: closure identity -> repaired source profile (Chapter 1).
% QGD-115 update, 26 October 2026: Einstein-Hilbert form derived by uniqueness (Lovelock) under principle L.
% QGD-113 final update, 26 October 2026: curved spin-1/2 flow law and slow-motion limits.
% QGD-112 update, 26 October 2026: exact spin-1/2 flow law; naming of derived objects corrected.
% QGD-111 update, 26 October 2026: covariant spin-1/2 version of the quantum terms.
% QGD-110 insertion, 26 October 2026: quantum structure of the master metric and
% covariance tests (from research notes QGD-107--109 and the QGD-110 tests).
% QGD-105 audit correction, 26 October 2026: QGD-101 insertions restricted to what
% follows from this chapter's stated assumptions; reframing of QGD's identity removed.
% All QGD-101 insertions are marked; no prior equation or label is removed.
% Units c=hbar=1 in variational arguments.
% QGD-136 integration, 31 October 2026: renormalized quantum-effective action, static collective matching, classical limit.
% QGD-124 insertion, 26 October 2026: abstract states QGD's place as quantum gravity (classical/emergent branch).
% QGD-123 housekeeping, 26 October 2026: naming convention applied to unambiguously derived objects;
% converged stellar radius (Ch2); benchmark reproducibility note (Ch3); spin-1/2 stress closed form (Ch4).
\documentclass[11pt,a4paper]{article}
\usepackage[T1]{fontenc}
\usepackage[utf8]{inputenc}
\usepackage[margin=25mm,headheight=15pt]{geometry}
\usepackage{amsmath,amssymb,amsthm,mathtools,bm,booktabs,microtype,xcolor,mathrsfs}
\usepackage[colorlinks=true,linkcolor=blue!55!black,urlcolor=blue!55!black]{hyperref}
\usepackage{fancyhdr}
\pagestyle{fancy}\fancyhf{}\lhead{Quantum Gravity Dynamics}\rhead{Chapter 4 — Natural-Invariant Dynamics}\cfoot{\thepage}
\newtheorem{theorem}{Theorem}[section]
\newtheorem{proposition}[theorem]{Proposition}
\newtheorem{remark}[theorem]{Remark}
\newtheorem{definition}[theorem]{Definition}
\newcommand{\dd}{\mathrm d}
\newcommand{\cE}{\mathcal E}
\newcommand{\cR}{\mathcal R}
\newcommand{\cG}{\mathscr G}
\newcommand{\cT}{\mathscr T}
\newcommand{\cQ}{\mathscr Q}
\newcommand{\cS}{\mathscr S}
\title{\textbf{Quantum Gravity Dynamics}\\[6pt]
\large Chapter 4: Natural-Invariant Dynamics\\ and the Physical Meaning of Gravity\\[5pt]
\normalsize QGD--098 research edition --- referee-verified radial Dirac audit and semiclassical benchmark}
\author{QGD research manuscript --- independent of the archived historical Chapter 4}
\date{5 November 2026 --- action-first edition; 31 October QGD--136 integration; QGD--101 integration, QGD--105 audit correction, QGD--110 quantum structure, QGD--111/112/113 spin-$\frac12$ updates, QGD--115 Lovelock}
\begin{document}
\maketitle
\begin{abstract}
What does the natural-invariant description of gravitation add to the
Einsteinian metric equations? We formulate a precise first answer, without
postulating a new wave law. Given the composite QGD metric of Chapters 1--2
and the effective gravitational and matter action used in Chapters 2--3,
variation with respect to the independent natural-invariant one-forms
produces an exact projected Euler equation. The historical three-channel
idea thereby admits an unambiguous \emph{action-defined successor},
$\cS=\cG+\cT+\cQ=0$. In the derivative-free, metric-only sector,
$\cQ=0$ and the remaining terms are projections of the Einstein--matter
residual. On a regular full-rank signed-field representation, those
projected equations are equivalent to the effective Einstein equations.
The result clarifies the physical role of $\sigma_\mu$, the origin of its
representation freedom, and the exact location where genuinely distinct
QGD dynamics would have to arise. We incorporate the Chapter~3 breakthroughs:
Newtonian calibration of the exterior natural-invariant amplitude; the
exact static Killing-energy expansion; a corrected conserved-current
bridge; the leading Dirac-derived quantum density-gradient stress; and a
no-go theorem for universal power-law exterior corrections produced solely
by restructuring a compact spherical source at fixed mass. It does not assert a microscopic
wave equation, a derived $\sigma$ kinetic term, or a new gravitational
observable. This short chapter is intended to grow as these questions
are solved; it does not alter the historical Chapter 4. The chapter proves a metric-only quantum-equivalence theorem, imports a
verified semiclassical massless-Dirac far-field correction, and audits the
historical $9/2$ force. A QGD--110 section derives the quantum structure of the
master metric from the founding principles: a universal quantum sector fixed by the
source's quantum stress, a per-wave quantum conformal factor, and the test of the
historical derivative term, and the exact flow law of spin-$\frac12$ matter in a curved
universal metric; it also records Lorentz-covariance and gauge tests of the main
equations. A QGD--115 subsection shows that the Einstein--Hilbert form of the
gravitational dynamics is not an independent axiom: it follows uniquely, by
Lovelock's theorem, from QGD's metric-only structure and a second-order locality
principle, with the coupling fixed by the Newtonian calibration. A QGD--120 subsection
states the complete QGD action, its variables, its symmetries and its quantum effective
form; a QGD--121 subsection shows that on spanning branches it propagates exactly the two
gravitational modes of general relativity. The revision history is in
Section~\ref{sec:revision-history}.
In the sense stated in Chapter~1's introduction, QGD is a theory of quantum gravity in the
classical/emergent branch: quantum matter sources a classical geometry, the quantum structure of
that geometry is derived here, and the metric itself is not quantized.
\end{abstract}

\section{Q061 final audit of the covariant QGD action and its extensions (5 November 2026)}\label{sec:q061-action-audit}
\subsection{Provenance: two axioms, one separate dynamical law}
The founding geometric postulate is the full $g_{\mu\nu}=T^a{}_{\mu}T^b{}_{\nu}[\eta_{ab}-\sum_s\varepsilon_s\sigma_a^{(s)}\sigma_b^{(s)}-\lambda\mathcal Q_{ab}]$. Axiom II postulates compatible microscopic realizability but fixes no classical or quantum action coefficient. The second-order, local, generally covariant \emph{metric-mediated} gravitational law used here is selected \emph{conditionally} on Principle~L; four-dimensional Lovelock uniqueness then gives the Einstein tensor and a cosmological constant, while Newtonian matching fixes $G_N$. This theorem does not exclude an additional direct-sector functional. The metric-only operational model always declares $\lambda\mathcal Q=0$ and restricts to a suitable regular field-map branch.

\subsection{Renormalization-conscious effective action}
For clarity, in units $\hbar=c=1$ except as a loop-order annotation, define a once-subtracted renormalized metric functional by
\begin{equation}\label{eq:q061-ch4-action}
\boxed{\begin{aligned}
\Gamma_{\rm met}^{\rm ren}[g,\psi_c,A_c;\varrho,\mu]
={}&\frac{1}{16\pi G_N}\int_{\mathcal M}\!d^4x\sqrt{-g}(R-2\Lambda_{\rm ren})+S_{\rm GHY}[g]
+S_{\rm F1}^{\rm cl}[\psi_c,A_c;g]\\
&+\int_{\mathcal M}\!d^4x\sqrt{-g}\left[c_R(\mu)R^2+c_C(\mu)C_{\alpha\beta\gamma\delta}C^{\alpha\beta\gamma\delta}
+c_E(\mu)E_4\right]\\
&+\Gamma_{\rm matter,rem}^{\rm ren}[g,\psi_c,A_c;\varrho,\mu]
+\Gamma_{\rm higher/bdy}.
\end{aligned}}
\end{equation}
Set $\Gamma_{\rm QGD}^{\rm ren}[\Sigma,\ldots]=\Gamma_{\rm met}^{\rm ren}[g(\Sigma),\ldots]$. The $c_i$ are independent finite-renormalized effective coefficients (including the matched quantum local parts) rather than guaranteed overall-order-$\hbar$ parameters. QN004 computes particular Euclidean one-loop form factors; the remainder in Eq.~\eqref{eq:q061-ch4-action} explicitly excludes local pieces already placed in $c_i$. $E_4$ is topological in four dimensions for fixed topology up to necessary boundary qualifications. Higher-derivative terms are treated perturbatively inside their EFT domain, not as a license to retain all formal runaway branches.

\subsection{Symmetries and the external-coframe qualification}
Under simultaneous transformations of all geometric fields, the metric-only action is spacetime-diffeomorphism covariant and local-frame Lorentz invariant; the Dirac spin connection and common electromagnetic $F_{\mu\nu}$ preserve spinor-frame covariance and local $U(1)$. For $Y$ the $N\times4$ sector matrix, $Y\mapsto O(x)Y$ with $O^{\mathsf T}CO=C$ leaves $g$ unchanged when $\lambda\mathcal Q=0$. These metric-preserving directions give no independent metric-only energy, even when $O(x)$ is nonconstant. In the declared variational scheme, however, the reference coframe $T$ is \emph{held fixed}. Formal covariance with $T$ transforming as a coframe does not itself establish unrestricted off-shell active diffeomorphism gauge freedom; the Noether identity must include any external reference background. An independently derivative-dependent $\mathcal Q$ requires its own covariant tensorial prescription and variational adjoint before it can be used.

\subsection{Causal quantum response and the additional-sector option}
For time-dependent expectation values use a Schwinger--Keldysh closed-time-path functional $\Gamma_{\rm CTP}[\Sigma_+,\Sigma_-;\varrho_{\rm in}]$, varying its forward history and then taking the physical limit $\Sigma_+=\Sigma_-$. Euclidean QN004 response functions cannot simply be substituted as retarded real-time kernels without an analytic/state prescription. A direct dark-energy functional has the separate form
\begin{equation}\label{eq:q061-ch4-extension}
\Gamma_{\rm ext}=\Gamma_{\rm QGD}^{\rm ren}+
\Delta S_{\rm sector}[\Sigma;\mathsf H_0,\text{parameters}],
\qquad
\frac{\delta\Gamma_{\rm ext}}{\delta\Sigma}
=(D_\Sigma g)^*\frac{\delta\Gamma_{\rm ext}}{\delta g}
+(D_\Sigma s)^*\frac{\delta\Gamma_{\rm ext}}{\delta s}
\end{equation}
when the new functional depends through $g(\Sigma)$ and an independent selected scalar $s(\Sigma)$. If the joint map $(g,s)$ has a right inverse on the admissible variation spaces, the full sector Euler equation is equivalent to the two separate metric and scalar Euler equations. The additional selector/reference, kinetic coefficient, and potential scale are new physical inputs. Signed sectors and universal metric coupling alone cannot determine them (Q055--Q060). Thus Eq.~\eqref{eq:q061-ch4-action}, not Eq.~\eqref{eq:q061-ch4-extension}, remains the canonical action.




\section*{The defining QGD action and the two-axiom hierarchy (5 November 2026)}
\addcontentsline{toc}{section}{The defining QGD action and two-axiom hierarchy}
\noindent\textbf{Action placement.} Chapter~1 states Axiom~I (complete master metric), Axiom~II (unproved independent microscopic realizability), and the separate defining variational law at the outset. This chapter is the central \emph{dynamical derivation}, not a deferred search for a law of gravity. The classical dynamics is the local Einstein--Hilbert--matter effective action expressed in independent natural-invariant variables, with renormalized matter loops added consistently:
\begin{equation}\label{eq:qgd-action-first-2026}
\boxed{\begin{aligned}
\Gamma_{\rm QGD}^{\rm ren}[\Sigma,\psi_c,A_c;\varrho,\mu]
={}&\frac{1}{16\pi G_N}\int_\mathcal M d^4x\sqrt{-g(\Sigma)}\,[R(g(\Sigma))-2\Lambda_{\rm ren}]\\
&+S_{\rm GHY}[g(\Sigma)]+S_{\rm F1}^{\rm cl}[\psi_c,A_c;g(\Sigma)]\\
&+\hbar\Gamma_{\rm m,ren}^{(1)}[g(\Sigma),\psi_c,A_c;\varrho,\mu]+O(\hbar^2).
\end{aligned}}
\end{equation}
Here $S_{\rm GHY}=(8\pi G_N)^{-1}\int_{\partial\mathcal M}\epsilon\sqrt{|h|}\,K$ with the standard boundary prescription, and
\begin{equation}
\boxed{S_{\rm F1}^{\rm cl}=\int d^4x\sqrt{-g}\left[\frac i2\bar\psi_c\gamma^\mu\overleftrightarrow\nabla_\mu\psi_c-m_0\bar\psi_c\psi_c+G_4(\bar\psi_c\psi_c)^2-\frac14F_{\mu\nu}F^{\mu\nu}\right].}
\end{equation}
The quantum term is decomposed in one renormalization scheme into local $\alpha_R R^2+\alpha_C C^2+\alpha_E E_4$, complementary nonlocal kernels, state contributions and mixed vertices. Chapter~6, Section~40, gives explicit one-loop Euclidean massive Dirac/Hartree form factors through curvature-quadratic order, not an arbitrary-curvature Lorentzian completion. These contributions must not be appended twice.

\noindent\textbf{Full-master equations and declared specialization.}
\begin{equation}\label{eq:qgd-optionC-first-2026}
\boxed{2s_s\Pi^{ab}\sigma_b^{(s)}=\lambda[(D_s\mathcal Q)^*_{\sqrt{-g}}\Pi]^a,\qquad
\Pi^{ab}=T^a{}_{\mu}T^b{}_{\nu}\frac1{\sqrt{-g}}\frac{\delta\Gamma_{\rm QGD}}{\delta g_{\mu\nu}}.}
\end{equation}
This is the action's chain-rule equation for fixed reference coframe and without additional direct $\Sigma$ dependence, provided $\mathcal Q$ is defined as a differentiable functional with its boundary prescription. In the \emph{defined} metric-only specialization $\lambda\mathcal Q=0$, it reduces to $\Pi^{ab}\sigma_b^{(s)}=0$. A full-rank metric Jacobian implies $\Pi^{ab}=0$, while rank-deficient sectors can be spuriously stationary. Axiom~II is not used to establish this result. If an independent direct sector action were added, its Euler derivative would modify this equation; no such addition is adopted here.

\noindent\textbf{Interpretive restriction.} Lovelock uniqueness motivates the local classical metric dynamics only after the hypotheses of locality, four dimensions and second-order metric Euler equations are imposed. It does not prove that this particular variational action follows from the geometric axioms alone. The present metric-only theory is physically equivalent to the corresponding metric EFT on regular complete branches unless additional \emph{specified and justified} sector dynamics or admissibility conditions yield distinct observables.


\noindent\textbf{Reading key.} \textbf{D}: derived from stated assumptions;
\textbf{A}: adopted effective dynamics; \textbf{C}: fixed definition or
normalization; \textbf{P}: proposed physical interpretation not yet an
independent empirical result. These labels distinguish logical status,
not scientific importance.

\section{From the quantum description to effective geometry}
\label{sec:motivation}
General relativity establishes how geometry behaves when coupled to
energy and momentum. It does not, by itself, establish whether geometry
is fundamental or an effective macroscopic description. QGD investigates
a concrete alternative interpretation: the gravitational field is
associated with macroscopic quantum structure of matter, represented by
gravitational natural-invariant covectors $\sigma_\mu^{(s)}$; their
composites define an effective metric (P). The word \emph{natural
invariant} is QGD's historical name for these covectors; a covector is
not itself a coordinate-invariant scalar.

The essential distinction is between an \emph{ontological} proposal and a
\emph{dynamical} result. The first three chapters supply (i) the
four-component macroscopic spinor structure and controlled local spinor
ratios, (ii) the master metric and its signed-sector representations,
and (iii) Einsteinian effective dynamics and Newtonian calibration in
specified limits. None of these statements alone proves that the
individual natural-invariant components constitute independent
observable particles or waves.

The complete founding QGD master metric is retained unchanged (the signed-sector-only formula is a specialization, not the master metric):
\begin{equation}\label{eq:master}
\boxed{g_{\mu\nu}=T^a{}_{\mu}T^b{}_{\nu}
\left[\eta_{ab}-\sum_{s=1}^{N}\varepsilon_s
\sigma_a^{(s)}\sigma_b^{(s)}-\lambda\mathcal Q_{ab}\right],
\quad \varepsilon_s=\pm1.}
\end{equation}
Here $T^a{}_{\mu}$ is the reference coframe, and the derivative
operator $\mathcal Q_{ab}$ and coefficient $\lambda$ remain
unmatched proposals in the unrestricted theory. We will derive definite
equations in the explicitly stated fixed-coframe, $\lambda=0$,
metric-only effective sector and then explain what changes when the
full metric map is used.

\section{The new Chapter 4 field equation}
\label{sec:balance}
In this section set $c=1$. Restore the physical Einstein coupling
$8\pi G_N/c^4$ when using SI units. Consider the composite effective
action (A)
\begin{equation}\label{eq:action}
\Gamma[\sigma,\Psi]=S_{\rm EH}[g(\sigma)]
+S_{\rm matter}[g(\sigma),\Psi],\qquad
S_{\rm EH}=\frac{1}{16\pi G_N}
\int\dd^4x\sqrt{-g}\,(R-2\Lambda).
\end{equation}
Boundary variations are compactly supported or supplied with the
appropriate gravitational boundary term. The matter equations are
varied independently. Define $T^{\mu\nu}$ by
$\delta S_{\rm matter}=\frac12\int\sqrt{-g}\,
T^{\mu\nu}\delta g_{\mu\nu}$ for on-shell matter variations, and let
\begin{equation}\label{eq:Eg}
\boxed{\cE_g^{\mu\nu}:=\frac1{\sqrt{-g}}
\frac{\delta\Gamma}{\delta g_{\mu\nu}}
=-\frac{1}{16\pi G_N}
\left(G^{\mu\nu}+\Lambda g^{\mu\nu}
-8\pi G_N T^{\mu\nu}\right).}
\end{equation}
This is a definition of the effective metric Euler derivative,
not an independent QGD field equation.

% ===================== QGD-115 INSERTION =====================
\subsection{Why the Einstein--Hilbert form: a uniqueness theorem (QGD--115)}
\label{sec:lovelock}
The action~\eqref{eq:action} was introduced as an adopted effective action (A).
Its \emph{form} need not be adopted. It follows from facts established within QGD
together with one structural principle.

\paragraph{Established within QGD.}
\begin{enumerate}
\item[(i)] \emph{Matter couples only through $g$.} The microscopic action uses a
vierbein of $g$ and is locally Lorentz invariant, so matter depends on the natural
invariants only through $g$ (Proposition~\ref{prop:nobeyond}).
\item[(ii)] \emph{Covariance.} The master metric is tensorial under coordinate
changes and invariant under frame Lorentz transformations
(Section~\ref{sec:quantumstructure}).
\item[(iii)] \emph{Conservation.} Diffeomorphism invariance of the matter action
gives $\nabla_\mu T^{\mu\nu}=0$ on shell.
\item[(iv)] \emph{No consistent dynamics beyond $g$ has been found.} The routes
tested, physical frame redundancies, an antisymmetric pairing field and two
dynamical frames, were either inconsistent at two-derivative order or required an
independent postulate (Section~\ref{sec:physics}).
\end{enumerate}

\paragraph{Principle L (locality and order).} The gravitational equation for the
natural invariants is local, generally covariant, depends on $\sigma$ only through
$g$, and contains at most second derivatives of $g$.

\begin{theorem}[Einstein form by uniqueness (D$\mid$L; Lovelock)]\label{thm:lovelock}
Under (i)--(iii) and Principle~L, in four spacetime dimensions the gravitational
equation has the form $\mathcal E_{\mu\nu}[g]=\kappa\,T_{\mu\nu}$ with
$\mathcal E_{\mu\nu}$ symmetric, built from $g$ and its first two derivatives, and
identically divergence free. By Lovelock's theorem \cite{Lovelock1971,Lovelock1972}
the only such tensors are $aG_{\mu\nu}+bg_{\mu\nu}$, so
\begin{equation}\label{eq:lovelock}
\boxed{G_{\mu\nu}+\Lambda g_{\mu\nu}=\kappa\,T_{\mu\nu}.}
\end{equation}
The Newtonian calibration of Chapter~3 then fixes $\kappa=8\pi G_N/c^4$, and the
action~\eqref{eq:action} is the corresponding variational form, up to boundary
terms.
\end{theorem}
\begin{proof}
The equation must be divergence free because its source is, by~(iii), for every
matter configuration. Lovelock's theorem classifies the symmetric, divergence-free
concomitants of $g$ with at most second derivatives in four dimensions. The
coupling is the single constant left, and F5 fixes it.
\end{proof}
\noindent\textbf{Verification.} On a generic metric with all ten components
varying, a 50-digit computation gives $\nabla^\mu G_{\mu\nu}=O(10^{-14})$, while
$\nabla^\mu R_{\mu\nu}$ and $\nabla^\mu(Rg_{\mu\nu})$ are of order $0.1$--$0.7$; the
unique divergence-free combination $aR_{\mu\nu}+bRg_{\mu\nu}$ has $b/a=-0.5000$.

\paragraph{Status.} The \emph{form} of the gravitational dynamics is derived
(D$\mid$L), and the \emph{coupling} is fixed by the calibration F5. Three things are
not derived:
\begin{itemize}
\item \textbf{$\Lambda$.} The theorem allows a cosmological term but does not fix its
value.
\item \textbf{The second-order requirement.} It is a principle, chosen for a
well-posed, ghost-free classical limit. Higher-derivative terms, such as the
logarithmic curvature-squared terms induced by quantum matter, are treated as
subleading corrections.
\item \textbf{The value of $G_N$ from the microscopic theory.} Two routes remain
open. In the induced-gravity route (QGD--114), the one-loop effective action of
Chapter~1's fields yields the Einstein form, but with one spin-$\frac12$ field and
one Maxwell field the induced Newton constant is negative and the induced
cosmological constant is of cutoff order \cite{Sakharov,Visser}. In the
thermodynamic route \cite{Jacobson}, horizon entropy counts physical degrees of
freedom positively and $\Lambda$ appears as an integration constant, but an
equilibrium hypothesis and a physical cutoff are required.
\end{itemize}
% =================== END QGD-115 INSERTION ===================

Write the metric-only signed specialization locally as
\begin{equation}\label{eq:signed}
g_{\mu\nu}=\bar g_{\mu\nu}
+\sum_s s_s\sigma^{(s)}_\mu\sigma^{(s)}_\nu,
\qquad s_s:=-\varepsilon_s\in\{+1,-1\}.
\end{equation}
All variations below treat the covectors as independent components
in a fixed frame and hold the background coframe fixed.

\begin{theorem}[Action-defined natural-invariant equation (D$\mid$A)]
\label{thm:projection}
For the action \eqref{eq:action} and metric map \eqref{eq:signed},
the complete Euler derivative of each independent natural-invariant
covector satisfies
\begin{equation}\label{eq:Euler}
\boxed{\cS_{(s)}^\rho
:=\frac1{\sqrt{-g}}\frac{\delta\Gamma}
{\delta\sigma^{(s)}_\rho}
=2s_s\cE_g^{\rho\nu}\sigma^{(s)}_\nu=0.}
\end{equation}
Equivalently, with \emph{action-defined} channels,
\begin{equation}\label{eq:split}
\boxed{\cS_{(s)}^\rho
=\underbrace{-\frac{s_s}{8\pi G_N}
\left(G^{\rho\nu}+\Lambda g^{\rho\nu}\right)
\sigma^{(s)}_\nu}_{\displaystyle\cG_{(s)}^\rho}
+\underbrace{s_s T^{\rho\nu}\sigma^{(s)}_\nu}
_{\displaystyle\cT_{(s)}^\rho}
+\underbrace{0}_{\displaystyle\cQ_{(s)}^\rho}=0.}
\end{equation}
\end{theorem}
\begin{proof}
Vary the signed metric while holding other sectors fixed:
$\delta g_{\mu\nu}=s_s(\delta\sigma^{(s)}_\mu
\sigma^{(s)}_\nu+\sigma^{(s)}_\mu\delta\sigma^{(s)}_\nu)$.
Symmetry of $\cE_g^{\mu\nu}$ gives
$\delta\Gamma=\int\sqrt{-g}\,2s_s\cE_g^{\rho\nu}
\sigma^{(s)}_\nu\delta\sigma^{(s)}_\rho$.
Equation \eqref{eq:Euler} follows from independent arbitrary
compactly supported variations. Substituting \eqref{eq:Eg}
proves \eqref{eq:split}.
\end{proof}

\begin{remark}[What the three symbols mean]
$\cG$ is the effective \emph{geometric-action contribution}, $\cT$
is the \emph{matter-stress contribution}, and $\cQ$ denotes additional
terms arising from a \emph{specified} derivative-dependent metric map
when such a map has actually been defined and varied. Thus $\cQ=0$
in the metric-only algebraic sector. The letter $Q$ does \emph{not}
by itself mean ``quantum''. In particular, a genuine quantum
matter-loop correction to a metric-only effective action modifies
$\cE_g$, hence the geometric/matter balance, without necessarily
producing any separate $\cQ$ channel.
\end{remark}


% ===================== QGD-120 INSERTION =====================
\subsection{The complete QGD action (QGD--120)}\label{sec:qgdaction}
The action~\eqref{eq:action} left $S_{\rm matter}$ generic. Every ingredient needed to
make it explicit has been derived in Chapters~1--4; here it is assembled.

\paragraph{Fields.} The dynamical fields are the natural invariants $\sigma^{(s)}_a$
($s=1,\dots,N$, signs $s_s=\pm1$), the microscopic spin-$\frac12$ field $\psi$ and the Maxwell
field $A_\mu$. The reference coframe $T^a{}_\mu$ is held fixed. The metric is the classical
restriction of the master metric,
\begin{equation}\label{eq:qgdmetric}
g_{\mu\nu}=T^a{}_\mu T^b{}_\nu\Bigl[\eta_{ab}+\sum_ss_s\sigma^{(s)}_a\sigma^{(s)}_b\Bigr].
\end{equation}
There is no independent derivative term $\lambda\mathcal Q$. The quantum content of the
master metric arises from the effective action below, not from an added term
(Section~\ref{sec:quantumstructure}).

\begin{definition}[QGD action]\label{def:qgdaction}
\begin{equation}\label{eq:qgdS}
\boxed{S_{\rm QGD}[\sigma,\psi,A]=\frac1{16\pi G_N}\int_{\mathcal M}\!\mathrm d^4x\sqrt{-g}\,(R-2\Lambda)
+\frac1{8\pi G_N}\oint_{\partial\mathcal M}\!\mathrm d^3y\,\epsilon\sqrt{|h|}\,K
+S_{\rm F1}[\psi,A;g]}
\end{equation}
with $g=g[\sigma;T]$ from~\eqref{eq:qgdmetric}, the Gibbons--Hawking--York term for a
well-posed variational principle, and the microscopic action of Chapter~1 written on $g$,
\begin{equation}\label{eq:SF1}
S_{\rm F1}=\int\!\mathrm d^4x\sqrt{-g}\Bigl[\tfrac i2\bigl(\bar\psi\gamma^\mu\nabla_\mu\psi-\overline{\nabla_\mu\psi}\gamma^\mu\psi\bigr)
-m_0\bar\psi\psi+G_4(\bar\psi\psi)^2-\tfrac14F_{\mu\nu}F^{\mu\nu}\Bigr].
\end{equation}
Here $G_4$ is Chapter~1's contact coupling (renamed to avoid a clash with $G_N$),
$\gamma^\mu=e_a{}^\mu\gamma^a$ with $e$ any vierbein of $g$,
$\{\gamma^\mu,\gamma^\nu\}=-2g^{\mu\nu}$, and
$\nabla_\mu=\partial_\mu+\Gamma_\mu+iqA_\mu$ with the spinor connection fixed by $\nabla\gamma=0$.
\end{definition}

\begin{theorem}[Symmetries of the QGD action (D)]\label{thm:qgdsym}
$S_{\rm QGD}$ is invariant under:
\begin{enumerate}
\item diffeomorphisms acting on all fields, $T$ included;
\item local Lorentz transformations of the frame index, $T\mapsto\Lambda T$,
$\sigma_a\mapsto\sigma_c(\Lambda^{-1})^c{}_a$, together with local Lorentz rotations of the vierbein
$e$ and the spinor $\psi$;
\item local sector rotations $\sigma\mapsto O\sigma$ with $O^{\mathsf T}CO=C$, $C=\operatorname{diag}(s_s)$;
\item local $U(1)$ gauge transformations of $(\psi,A)$.
\end{enumerate}
Any configuration with $\sum_ss_s\sigma^{(s)}_a\sigma^{(s)}_b=0$ and flat $T$ gives Minkowski space,
whose isometry group is the Poincar\'e group. For perturbation theory the vacuum must be
represented on a spanning branch, not at $\sigma=0$ (Theorem~\ref{thm:dof}).
\end{theorem}
\begin{proof}
The gravitational terms depend on $(\sigma,T)$ only through $g$, which is invariant under
2 and 3 and a tensor under 1 (Section~\ref{sec:quantumstructure}, covariance table). The
matter term depends on $g$ through a vierbein; its local Lorentz invariance follows from
$\nabla\gamma=0$, under which $[\nabla_\mu,\nabla_\nu]=-\frac14R_{ab\mu\nu}\gamma^a\gamma^b$ (verified in
Chapter~4, Theorem~\ref{thm:spinhalf}). The $U(1)$ invariance is standard.
\end{proof}

\paragraph{Variables and the role of $T$.} $S_{\rm QGD}$ is varied with respect to
$\sigma$, $\psi$ and $A$; $T$ is not varied. This does not break covariance in effect. On
any branch where the metric map has a right inverse, every metric variation is produced by
a variation of $\sigma$ (Chapter~2, Correspondence Theorem), so a change of $T$ can be
compensated by $\sigma$, and the physical content is independent of the choice of $T$.
The individual $\sigma^{(s)}$ are therefore not gauge-invariant observables: only $g$ and
quantities built from it are (Section~\ref{sec:quantumstructure}).

\paragraph{Field equations.} Varying $\sigma$ gives the field-balance
equation~\eqref{eq:Euler}. On spanning branches it is equivalent to
$G_{\mu\nu}+\Lambda g_{\mu\nu}=8\pi G_NT_{\mu\nu}$, whose form is unique by
Theorem~\ref{thm:lovelock}. Varying $\psi$ gives the macroscopic-limit spinor equation on $g$;
varying $A$ gives Maxwell's equations with the $U(1)$ current.

\noindent\textbf{Verification.} For the static branch
$\sigma=A(r)\,\mathrm dt+B(r)\,\mathrm dr$ on the flat spherical reference, the Euler--Lagrange
equations of $\sqrt{-g}\,R$ computed directly in sector variables are exactly twice those of
the reduced Lagrangian $(rF)'/\sqrt D+\sqrt D$ used in Chapters~2--3, for both $A$ and $B$.
That reduced Lagrangian is therefore the Einstein--Hilbert action in natural-invariant
variables, up to normalization and a total derivative.

\paragraph{Quantum effective action.} At the quantum level the matter fields are
integrated out on $g(\sigma)$:
\begin{equation}\label{eq:GammaQGD}
\begin{aligned}
\Gamma_{\rm QGD}[\sigma]&=\frac{1}{16\pi G_N}\int\!\sqrt{-g}\,(R-2\Lambda)
+S_{\rm GHY}+\Gamma_{\rm F1}[g(\sigma)],\\
e^{i\Gamma_{\rm F1}[g]/\hbar}&=\int\mathcal D\bar\psi\,\mathcal D\psi\,\mathcal DA\;e^{iS_{\rm F1}[\psi,A;g]/\hbar}.
\end{aligned}
\end{equation}
with gauge fixing and a regulator; F2's coarse-graining is the Wilsonian version of
$\Gamma_{\rm F1}$. Stationarity of $\Gamma_{\rm QGD}$ gives the semiclassical equations with
$\langle T_{\mu\nu}\rangle$. The quantum structure derived in Section~\ref{sec:quantumstructure}, namely the
universal sector $q$, the per-wave factor $\Omega_\psi$, and the spin-$\frac12$ flow law, consists of
consequences of~\eqref{eq:GammaQGD}, not additional terms.

\paragraph{Status.}
\begin{center}\small
\begin{tabular}{@{}p{0.45\linewidth}p{0.47\linewidth}@{}}\toprule
Element & Status\\\midrule
Einstein--Hilbert form & Derived under Principle~L (Theorem~\ref{thm:lovelock})\\
$G_N$ & Fixed by the Newtonian calibration (Chapter~3)\\
$\Lambda$ & Allowed, not determined\\
Matter sector $S_{\rm F1}$ & Chapter~1's microscopic action, covariantly coupled\\
Coupling of matter to $\sigma$ & Only through $g$ (Proposition~\ref{prop:nobeyond})\\
Derivative term $\lambda\mathcal Q$ & Absent; the historical trial is excluded (Proposition~\ref{prop:trialtest})\\
Symmetries & Diffeomorphisms, local Lorentz, local sector rotations, $U(1)$\\
Quantum level & $\Gamma_{\rm QGD}$: matter quantized on a classical $g(\sigma)$\\
Quantization of $\sigma$ itself & Not defined: it would require a measure for $\sigma$ and gauge fixing
of the sector redundancy, and is not needed under the semiclassical reading\\\bottomrule
\end{tabular}
\end{center}
\noindent In summary, $S_{\rm QGD}$ is a well-defined, generally covariant and locally
Lorentz-invariant classical action, and $\Gamma_{\rm QGD}$ is its well-defined semiclassical
quantum effective action. A path integral over the natural invariants themselves is not
part of QGD as formulated.
% =================== END QGD-120 INSERTION ===================

% ====================== QGD-136 RESEARCH COMPLETION ======================
\subsection{Renormalized final action and quantum-to-classical limit (QGD--136)}
\label{sec:qgd136final}
\noindent\textbf{Scope and status.}
This subsection consolidates the Chapter~4 action with the conditional static
matter-wave correspondence of Chapter~5 v3.0 and the massless-Dirac
semiclassical correction of Chapter~5 v3.1. It is a proposed \emph{effective,
semiclassical} final form, not a derivation of an autonomous collective
metric from Chapter~1's flat-space microscopic theory. We work with $c=1$
and display $\hbar$ as a formal loop-counting parameter. To discuss a
classical limit, hold fixed the \emph{renormalized} $G_N$, the macroscopic
classical matter source and appropriate boundary data; a naive $\hbar\to0$
limit of an arbitrarily normalized microscopic spinor is not the same limit.

\begin{definition}[Quantum-completed QGD effective action (EFT)]\label{def:qgd136}
For the explicitly controlled and currently accepted semiclassical branch,
set $\lambda\mathcal Q_{ab}=0$ in the as-yet-unmatched historical
metric-derivative proposal, and take the reduced classical metric map
\begin{equation}\label{eq:qgd136metric}
 g_{\mu\nu}[\sigma;T]=T^a{}_{\mu}T^b{}_{\nu}
 \left[\eta_{ab}+\sum_s s_s\sigma_a^{(s)}\sigma_b^{(s)}\right].
\end{equation}
On a regular signed-sector branch, define the \emph{renormalized}
semiclassical action, with a specified quantum matter state $\varrho$,
\begin{equation}\label{eq:qgd136final}
\boxed{\begin{aligned}
\Gamma_{\rm QGD}^{\rm ren}[\sigma,\psi_c,A_c;\varrho]={}&
\frac{1}{16\pi G_N}\int_{\mathcal M}\!\dd^4x\sqrt{-g(\sigma)}
\,[R(g(\sigma))-2\Lambda_{\rm ren}]\\
&+S_{\rm GHY}[g(\sigma)]
+S_{\rm F1}^{\rm cl}[\psi_c,A_c;g(\sigma)]\\
&+\hbar\,\Gamma_{\rm m,ren}^{(1)}
[g(\sigma),\psi_c,A_c;\varrho;\mu]
+O(\hbar^2).
\end{aligned}}
\end{equation}
Here $\psi_c,A_c$ are mean/classical matter backgrounds where retained.
The quantum functional includes one-loop effects of fluctuating matter
fields \emph{not already present in the classical background action}.
The matter action $S_{\rm F1}^{\rm cl}$ is precisely Eq.~\eqref{eq:SF1},
with its spin connection, Maxwell and four-fermion components. When
\emph{all} matter is integrated out rather than retaining $\psi_c,A_c$,
one uses the equivalent background-free quantum influence functional and
must not also count the same matter stress in $S_{\rm F1}^{\rm cl}$.
The finite one-loop functional is defined using a consistent regulator,
subtraction scheme, renormalization scale and appropriate boundary/state
prescription. A finite-cutoff interpretation is allowed for Chapter~1's
four-fermion effective interaction.
\end{definition}

\paragraph{What the symbol $\Gamma_{\rm m,ren}^{(1)}$ contains.}
The quantum term of Eq.~\eqref{eq:qgd136final} is \emph{not} merely a
chosen local correction to the natural-invariant metric. In four dimensions
its renormalized curvature-dependent part has schematically
\begin{equation}\label{eq:qgd136loops}
\begin{aligned}
\Gamma_{\rm m,ren}^{(1)}={}&
\int\!\dd^4x\sqrt{-g}\,
\bigl[a_R(\mu)R^2+a_C(\mu)C_{\mu\nu\rho\sigma}C^{\mu\nu\rho\sigma}\bigr]
\\&+\Gamma_{\rm nonlocal}^{(1)}[g;\mu]
+\Gamma_{\rm state}^{(1)}[g;\varrho]
+\Gamma_{\rm vertices}^{(1)}[g,\psi_c,A_c;\varrho],
\end{aligned}
\end{equation}
with local curvature-squared renormalization and, for example,
nonlocal curvature kernels of the form
$\int\sqrt{-g}\,C\log(-\Box/\mu^2)C$ and
$\int\sqrt{-g}\,R\log(-\Box/\mu^2)R$ when justified by the
relevant massless-field calculation. The coefficients and admissible
kernels are theory-, state- and scheme-dependent and are \emph{not}
freely fixed to reproduce a desired force. Topological and boundary
counterterms may be needed and are understood in the chosen variational
problem. Equations~\eqref{eq:qgd136final}--\eqref{eq:qgd136loops} are an
organizing form, \emph{not} a claim that Chapters~1--5 have calculated
every one-loop form factor. For time-dependent semiclassical backreaction
the appropriate object is a causal in--in (Schwinger--Keldysh) effective
action; an in--out determinant is not generally a causal evolution law.

\begin{proposition}[Controlled Einstein limit, D$\mid$A]\label{prop:qgd136limit}
Suppose (i) $\lambda\mathcal Q$ is absent or tends to zero,
(ii) any independent classical $R^2$, $C^2$ and other higher-curvature
couplings are set to zero at the selected classical matching point,
(iii) the renormalized $G_N$ and $\Lambda$ remain finite in the limit,
(iv) the matter state admits a smooth classical macroscopic stress
limit, and (v) all loop/higher-order remainders in
Eq.~\eqref{eq:qgd136final} vanish as $\hbar\to0$ on the specified
regular branch. Then the full-rank QGD natural-invariant variation
reduces to the Einstein--matter equations,
\begin{equation}\label{eq:qgd136limit}
\boxed{\lim_{\hbar\to0}\frac{\delta\Gamma_{\rm QGD}^{\rm ren}}
 {\delta\sigma^{(s)}_\rho}=0
\quad\Longleftrightarrow\quad
G_{\mu\nu}+\Lambda g_{\mu\nu}=8\pi G_NT_{\mu\nu}^{\rm cl}.}
\end{equation}
\end{proposition}
\begin{proof}
The quantum term and higher-order remainders vanish by hypotheses
(iv)--(v). The remaining action is the Einstein--Hilbert--GHY plus
classical matter action at fixed renormalized couplings. Its metric
Euler tensor is $\mathcal E_g^{\mu\nu}\propto
G^{\mu\nu}+\Lambda g^{\mu\nu}-8\pi G_NT_{\rm cl}^{\mu\nu}$.
The composite metric variation is
$\delta g_{\mu\nu}=2\sum_ss_s\sigma_{(\mu}^{(s)}
\delta\sigma_{\nu)}^{(s)}$. Hence stationarity in the sector
variables gives $2s_s\mathcal E_g^{\rho\nu}\sigma_\nu^{(s)}=0$.
Surjectivity of the metric map on a regular signed-sector branch makes
this equivalent to $\mathcal E_g^{\mu\nu}=0$. The claimed classical
Einstein equations follow.
\end{proof}

\paragraph{Where Chapter~5's $g_c^2/(4\pi Z)$ enters (not an extra kinetic term).}
The established \emph{conditional} static collective-field model of
Chapter~5 v3.0 uses $\mathcal E[q_{\rm coll};\rho_m]
=\int\dd^3x\,[\tfrac Z2|\bm\nabla q_{\rm coll}|^2-g_cq_{\rm coll}\rho_m]$,
$Z>0$. Its result
\begin{equation}\label{eq:qgd136match}
\boxed{G_{\rm eff}=\frac{g_c^2}{4\pi Z},\qquad
G_{\rm eff}\stackrel{\rm calibration}{=}G_N,\qquad
H=\frac{2g_cq_{\rm coll}}{c^2}=-\frac{2\Phi_N}{c^2}}
\end{equation}
provides a \emph{static infrared matching condition} on the coefficient
already appearing in the Einstein--Hilbert term. Under this change of
variables the collective energy becomes
\begin{equation}\label{eq:qgd136static}
\boxed{\mathcal E_H=\int\!\dd^3x\left[
\frac{c^4}{32\pi G_{\rm eff}}|\bm\nabla H|^2
-\frac{c^2}{2}\rho_mH\right],\qquad
\nabla^2H=-\frac{8\pi G_{\rm eff}}{c^2}\rho_m.}
\end{equation}
Here $q_{\rm coll}$ denotes Chapter~5's collective scalar coordinate, not Chapter~4's distinct quantum metric deformation $q_{ab}$. This equality is exact \emph{inside the stated static scalar effective
model}; it is not a derivation of its mode from the original microscopic
Dirac action or of the full off-shell Einstein--Hilbert functional.
If $q_{\rm coll}$ is identified with the already-represented Newtonian gravitational
mode, its energy must \emph{not} also be appended to
Eq.~\eqref{eq:qgd136final}, as this would double count the same physical
static response. A separate dynamical $q_{\rm coll}$ field can be included only
if its additional physical degrees of freedom and interactions are
independently established.

\paragraph{What the $\hbar$-dependent force means.}
The massless-Dirac Boulware far-field result of
Proposition~\ref{prop:diracfarfield} is obtained by varying the
renormalized quantum matter functional on a classical effective metric,
and solving the resulting semiclassical equations. At its stated
asymptotic order, the supported-static-body Killing-energy potential
has quantum term
\begin{equation}\label{eq:qgd136force}
\boxed{\delta\Phi_Q(r)=-\frac{2G_N^2M\hbar}{15\pi c^3r^3},\qquad
\delta F_r^{\rm coordinate}=-m_t\,\partial_r\delta\Phi_Q
=-\frac{2G_N^2Mm_t\hbar}{5\pi c^3r^4}.}
\end{equation}
Here $F_r^{\rm coordinate}$ denotes the radial derivative of this
specified Killing-energy potential; it is not automatically the
locally measured proper holding force. The Boulware state,
asymptotically flat source, massless four-component Dirac field and
far-field limit are essential. This is the \emph{same} one-loop quantum
matter effect already represented in the metric solution, not a
second independent $-\lambda\mathcal Q$ contribution. The two
reconstructed exterior natural-invariant components need not be
null at this quantum order (Eq.~\eqref{eq:scsigmaAB}). It is not
a QGD-only prediction beyond semiclassical Einstein gravity.

\begin{remark}[Scope of the word ``final'']
Eq.~\eqref{eq:qgd136final} is the most complete \emph{currently
justified semiclassical effective action}, not a UV-complete
quantum-gravity action. The signed-sector variables themselves are
not quantized here. A genuine QGD-specific correction would require
a newly justified independent functional $\Delta\Gamma_{\rm sector}$;
we set $\Delta\Gamma_{\rm sector}=0$ in the accepted action,
rather than include an unspecified operator as if it were derived.
The effective action's metric-only character preserves Chapter~4's
quantum-equivalence result: the same renormalized matter content and
state yield the same observables as semiclassical GR in the regime.
\end{remark}
% =================== END QGD-136 INSERTION =========================



% ================= QGD-137 FULL MASTER METRIC VARIATION ===================
\section{Full master-metric variation and the QGD equivalent of Einstein's equations (QGD--137)}
\label{sec:qgd137-fullvary}
\textbf{Scope and correction.} The signed-sector identity
$g=\bar g+\sum_s s_s\sigma^{(s)}\otimes\sigma^{(s)}$
is \emph{only} the fixed-reference, $\lambda\mathcal Q=0$ specialization
of the founding master metric. It must never be promoted to the unrestricted
QGD definition. In this section the reference coframe $T^a{}_{\mu}$ is fixed
as a variational background (but transforms tensorially under a coordinate
change), and the covectors are independent variables. The complete founding
metric reads
\begin{equation}\label{eq:qgd137-master}
\boxed{g_{\mu\nu}=T^a{}_{\mu}T^b{}_{\nu}
\mathcal M_{ab},\quad
\mathcal M_{ab}=\eta_{ab}-\sum_{s=1}^N\varepsilon_s
\sigma_a^{(s)}\sigma_b^{(s)}-\lambda\mathcal Q_{ab},
\quad \varepsilon_s=\pm1.}
\end{equation}
The historical $\mathcal Q_{ab}$ is a \emph{proposed} symmetric
covariant correction, not a unique operator derived in Chapters~1--5.
The different source-induced metric deformation $q_{ab}$ discussed in
Section~\ref{sec:quantumstructure} is a \emph{solution-level} quantum
backreaction defined through the semiclassical field equations. It is
not to be inserted as a second independently varied term when the
same renormalized matter functional already generates its effect.

\subsection{The metric Euler tensor including quantum matter}
For the renormalized metric-only action of Eq.~\eqref{eq:qgd136final},
with its metric argument now interpreted as the full map
\eqref{eq:qgd137-master} \emph{when that map has been mathematically
specified}, write
\begin{equation}\label{eq:qgd137-Pi}
\mathcal E_g^{\mu\nu}:=\frac{1}{\sqrt{-g}}
\frac{\delta\Gamma_{\rm metric}^{\rm ren}}{\delta g_{\mu\nu}},
\qquad
\Pi^{ab}:=T^a{}_{\mu}T^b{}_{\nu}\mathcal E_g^{\mu\nu}.
\end{equation}
Use the Chapter~4 stress convention
$\delta S_{\rm matter}=(1/2)\int\sqrt{-g}\,T^{\mu\nu}\delta g_{\mu\nu}$,
and define the renormalized quantum source by
\begin{equation}\label{eq:qgd137-Tq}
T_{\rm q}^{\mu\nu}:=\frac{2}{\sqrt{-g}}\,
\frac{\delta(\hbar\Gamma_{\rm m,ren}^{(1)}+O(\hbar^2))}
{\delta g_{\mu\nu}}.
\end{equation}
Within the chosen single-metric variational shorthand,
\begin{equation}\label{eq:qgd137-Eg}
\boxed{\mathcal E_g^{\mu\nu}=
-\frac{1}{16\pi G_N}\left[
G^{\mu\nu}[g]+\Lambda_{\rm ren}g^{\mu\nu}
-8\pi G_N\bigl(T_{\rm cl}^{\mu\nu}+T_{\rm q}^{\mu\nu}\bigr)
\right].}
\end{equation}
The $T_{\rm q}$ notation covers matter-loop/state backreaction;
for causal real-time evolution one must extract the physical equations
from an in--in (closed-time-path) functional, not identify an
in--out effective action with a causal time evolution equation.

\subsection{The exact projected natural-invariant Euler equation}
\begin{theorem}[Full master-map stationarity (D$\mid$metric-only action)]
\label{thm:qgd137-projection}
Let $\mathcal Q_{ab}=\mathcal Q_{ab}[\sigma;T]$ be a specified,
functionally differentiable symmetric operator built from the
covectors and fixed reference data; assume the metric is Lorentzian
and boundary terms make the variation legitimate. Write
$J_\sigma:=D_\sigma g$ for the \emph{complete} metric-map derivative.
Stationarity in every independent sector is exactly
\begin{equation}\label{eq:qgd137-adjoint}
\boxed{J_\sigma^*\bigl(\sqrt{-g}\,\mathcal E_g\bigr)=0,}
\end{equation}
or, in reference-frame notation with the volume-weighted adjoint,
\begin{equation}\label{eq:qgd137-sector}
\boxed{\mathscr E_{(s)}^{\,c}:=\frac{1}{\sqrt{-g}}
\frac{\delta\Gamma_{\rm metric}^{\rm ren}}
{\delta\sigma_c^{(s)}}
=-2\varepsilon_s\Pi^{cb}\sigma_b^{(s)}
-\lambda\left[(D_s\mathcal Q)^*_{\sqrt{-g}}\Pi\right]^c=0.}
\end{equation}
These are the \emph{full metric-pullback field equations} of QGD.
If a separately justified $\Delta\Gamma_{\rm sector}[\sigma]$
exists, its independent Euler derivative must be added to the
left-hand side and the metric-only equivalence below need not hold.
\end{theorem}
\begin{proof}
With $\delta T^a{}_{\mu}=0$,
\begin{equation}\label{eq:qgd137-dg}
\delta g_{\mu\nu}=T^a{}_{\mu}T^b{}_{\nu}
\left[-\sum_s\varepsilon_s\bigl(
\delta\sigma_a^{(s)}\sigma_b^{(s)}+
\sigma_a^{(s)}\delta\sigma_b^{(s)}\bigr)
-\lambda\delta\mathcal Q_{ab}\right].
\end{equation}
Substitute into
$\delta\Gamma_{\rm metric}^{\rm ren}=\int\sqrt{-g}\,
\mathcal E_g^{\mu\nu}\delta g_{\mu\nu}$.
Symmetry of $\Pi^{ab}$ yields the algebraic contribution
$-2\varepsilon_s\Pi^{cb}\sigma_b^{(s)}$; integrate the
$D_s\mathcal Q$ contribution by parts, with admissible
variations/boundary conditions, to obtain its formal adjoint.
This proves Eqs.~\eqref{eq:qgd137-adjoint}--\eqref{eq:qgd137-sector}.
\end{proof}

For a local trial depending only on sectors and their first
reference derivatives, write
$\mathcal Q_{ab}=\mathcal Q_{ab}(\sigma,\partial\sigma;T)$.
Define
$C_{ab}^{(s)c}:=\partial\mathcal Q_{ab}/\partial\sigma_c^{(s)}$
and
$D_{ab}^{(s)\mu c}:=\partial\mathcal Q_{ab}/\partial(\partial_\mu\sigma_c^{(s)})$.
Then the explicit (coordinate) Euler equation is
\begin{equation}\label{eq:qgd137-expanded}
\boxed{\begin{aligned}
0={}&-2\varepsilon_s\Pi^{cb}\sigma_b^{(s)}
-\lambda\Pi^{ab}C_{ab}^{(s)c}\\
&+\frac{\lambda}{\sqrt{-g}}\partial_\mu
\left(\sqrt{-g}\,\Pi^{ab}D_{ab}^{(s)\mu c}\right).
\end{aligned}}
\end{equation}
Derivatives of $\mathcal Q$ and the volume density are included,
not silently set to zero. If instead $\mathcal Q$ uses the
\emph{dynamical} Levi--Civita derivative of $g$, then
$D_g\mathcal Q$ also contributes to $\delta\mathcal Q$;
Eq.~\eqref{eq:qgd137-expanded} cannot be used unchanged.
The full metric may then be implicit, and the appropriate full
$J_\sigma$ must first be determined.

\subsection{When the projected equations are equivalent to EFE}
\begin{theorem}[Conditional quantum-corrected Einstein correspondence]
\label{thm:qgd137-efe}
Assume (i) the physical action depends on the covectors only through
$g[\sigma]$, (ii) the full linearized metric map $J_\sigma$ has a
right inverse on the admissible symmetric metric variations with the
chosen boundary conditions, and (iii) the renormalized matter
variations and their state are treated consistently. Then
\begin{equation}\label{eq:qgd137-efe}
\boxed{\mathscr E_{(s)}^{\,c}=0\quad\forall s,c
\quad\Longleftrightarrow\quad
G_{\mu\nu}[g]+\Lambda_{\rm ren}g_{\mu\nu}
=8\pi G_N\bigl(T^{\rm cl}_{\mu\nu}+T^{\rm q}_{\mu\nu}\bigr).}
\end{equation}
In particular the equivalence is exact on Chapter~2's regular
full-rank signed-sector branch at $\lambda=0$. This theorem does
not assert that the right-inverse property holds for an arbitrary,
nonzero historical $\lambda\mathcal Q$.
\end{theorem}
\begin{proof}
If $\mathcal E_g=0$, then $J_\sigma^*\mathcal E_g=0$ by
Eq.~\eqref{eq:qgd137-adjoint}. Conversely, for any admissible
$\delta g$, choose $\delta\sigma$ with
$J_\sigma\delta\sigma=\delta g$. Vanishing of the sector
variation implies $\int\sqrt{-g}\mathcal E_g^{\mu\nu}\delta g_{\mu\nu}=0$
for every symmetric variation, hence $\mathcal E_g=0$ and
Eq.~\eqref{eq:qgd137-efe} follows from Eq.~\eqref{eq:qgd137-Eg}.
\end{proof}

\begin{proposition}[Stability of the full-rank branch under a controlled correction]
\label{prop:qgd137-smalllambda}
Let the linearized map act between fixed Banach spaces of admissible
sector and metric perturbations. Suppose the $\lambda=0$
map $J_0$ admits a \emph{bounded} right inverse $R_0$,
$J_0R_0=I$, and $J_\lambda=J_0+\lambda B$ on these spaces,
where $BR_0$ is bounded and $\|\lambda BR_0\|<1$.
Then $J_\lambda$ also has a bounded right inverse, namely
\begin{equation}\label{eq:qgd137-rightinverse}
\boxed{R_\lambda=R_0\left(I+\lambda BR_0\right)^{-1},
\qquad J_\lambda R_\lambda=I.}
\end{equation}
Consequently the full-map EFE correspondence survives every
\emph{specified} correction satisfying these regularity and
smallness conditions. These conditions are not verified for an
unspecified historical $\mathcal Q$.
\end{proposition}
\begin{proof}
$J_\lambda R_0=I+\lambda BR_0$; the Neumann series converges
for $\|\lambda BR_0\|<1$. Multiplication by its inverse gives
Eq.~\eqref{eq:qgd137-rightinverse}.
\end{proof}

\subsection{Classical limit, the quantum metric term, and the collective bridge}
If the chosen completion has $\lambda\mathcal Q\to0$ as
$\hbar\to0$, loop terms vanish at fixed measured $G_N,\Lambda$
and smooth classical sources, and the metric Jacobian remains
regular, Eq.~\eqref{eq:qgd137-efe} reduces to
$G_{\mu\nu}+\Lambda g_{\mu\nu}=8\pi G_NT^{\rm cl}_{\mu\nu}$.
For finite $\hbar$, the same metric-only action gives the
quantum-stress source in Eq.~\eqref{eq:qgd137-efe}; solving for
$g$ and decomposing it afterward as $g=g_{\rm cl}+q$ does not
require adding a second independent $q$ kinetic term.
Chapter~5's $G_N=g_c^2/(4\pi Z)$ remains a \emph{static matching
condition} on the measured coupling, not a new metric operator.
The per-wave $\widetilde g^{(\psi)}=\Omega_\psi^2g$
of Section~\ref{sec:quantumstructure} remains an effective
phase-trajectory description for a specified probe, not the
universal spacetime metric inserted into the gravitational action.
If the reference coframe is itself promoted to a dynamical
field, its variation yields \emph{additional} Euler equations;
this section does not silently vary $T$.
% ================= END QGD-137 FULL MASTER METRIC VARIATION ===================

\subsection{Degrees of freedom of the QGD action (QGD--121)}\label{sec:dof}
$S_{\rm QGD}$ is written in terms of $4N$ sector components, more than the ten components of
$g$. Around a background $\bar\sigma$, the action depends on $\delta\sigma$ only through
$\delta g=D(\bar\sigma)\,\delta\sigma$, so its quadratic kernel is
\begin{equation}\label{eq:Ksigma}
K_\sigma(k)=D^{\mathsf T}K_{\rm FP}(k)\,D,
\end{equation}
with $K_{\rm FP}$ the linearized Einstein--Hilbert kernel. Gauge directions are the null
vectors of $K_\sigma$ at generic $k$. Propagating modes are the additional null vectors that
appear on the light cone, and dynamical components are those with a $k_0^2$ term.

\begin{theorem}[Degrees of freedom (D, N)]\label{thm:dof}
On any spanning branch ($\operatorname{rank}D=10$):
\begin{enumerate}
\item the $4N-10$ directions of $\ker D$ are the tangent space of the local sector-rotation
orbit ($O(1,3)$ for four $\eta$-signed sectors, $O(4,4)$ for eight). Since this equals the
dimension of the fibre of $\sigma\mapsto g$, the fibres are locally orbits, and the
redundancy is an exact local gauge symmetry of the nonlinear action;
\item together with the four diffeomorphism directions these exhaust the gauge null
space; the $k_0^2$ part of $K_\sigma$ has rank $6$, as in general relativity;
\item exactly \textbf{two} modes propagate: $4N-(4N-10)-4-4=2$.
\end{enumerate}
\end{theorem}
\begin{center}\small
\begin{tabular}{@{}lcccccc@{}}\toprule
Background & comps & $\ker D$ & null, generic $k$ & null, $k^2=0$ & propagating & $k_0^2$ rank\\\midrule
GR (reference) & 10 & -- & 4 & 6 & 2 & 6\\
4 sectors, $\sigma^{(a)}=v\,\mathrm dx^a$ & 16 & 6 & 10 & 12 & 2 & 6\\
8 sectors, cancelling background & 32 & 22 & 26 & 28 & 2 & 6\\
8 sectors, random background & 32 & 22 & 26 & 28 & 2 & 6\\
$\sigma=0$ & 16 & 16 & 16 & 16 & 0 & 0\\
1 sector, timelike & 4 & 0 & 2 & 2 & 0 & 0\\
1 sector, null & 4 & 0 & 2 & 2 & 0 & 2\\\bottomrule
\end{tabular}
\end{center}
\noindent\textbf{Verification.} Ranks computed numerically from Eq.~\eqref{eq:Ksigma}; GR's
own count ($4$, $6$, $2$) reproduced as a control. At random spanning backgrounds the
group orbit has dimension $6$ (four sectors) and $22$ (eight sectors), equal to $\dim\ker D$,
with every orbit direction in the kernel to $10^{-16}$. Random backgrounds and random momenta
always give two propagating modes.

\begin{proposition}[Degenerate and non-spanning branches (D, N)]\label{prop:degenerate}
\begin{enumerate}
\item At $\sigma=0$, $D=0$ and the quadratic action vanishes identically. The metric there
is Minkowski, but $\sigma=0$ is a degenerate point of the variables: perturbation theory
about it has no kinetic terms. The vacuum must be represented on a spanning branch, such as
$\sigma^{(a)}=v\,\mathrm dx^a$ with signs $(-1,+1,+1,+1)$ or the eight-sector cancelling background,
which describe the same metric with non-degenerate variables.
\item On one-sector branches nothing propagates at linear order: the timelike branch has no
time-derivative terms, and the null (Kerr--Schild-type) branch has two but no light-cone
modes. The sector equations there are projections of Einstein's equations (Chapter~2) and do
not reproduce the radiative sector of general relativity, consistent with Chapter~2's
pp-wave counterexample.
\end{enumerate}
\end{proposition}
\noindent\textbf{Consequence.} On spanning branches the extra sector components are pure
gauge, there are no additional or ghost modes, and the QGD action has exactly the dynamical
content of general relativity. The spanning condition, previously a requirement for the
equivalence of the field equations, is also the condition for the action to have the right
degrees of freedom.

\begin{proposition}[Spanning representations of the chapters' backgrounds (QGD--122)]\label{prop:backgrounds}
\begin{enumerate}
\item \textbf{Kerr--Schild class.} For any metric $g=\eta+H\,k\otimes k$ with $k$ null
(Schwarzschild, Reissner--Nordstr\"om, Vaidya, the regular cores, the pp-waves, Kerr and
Kerr--Newman), the four-sector lift
$\Sigma=v\,\mathbb 1+\frac H{2v}(\eta k)\,k^{\mathsf T}$ with signs $(-1,+1,+1,+1)$ and flat reference
$(1-v^2)\eta$ ($0<v<1$) reproduces $g$ exactly, and $\det\Sigma=v^4$ because $k$ is null. It is
therefore spanning wherever $H$ is finite, including at and inside horizons.
\item \textbf{Cosmology.} Chapter~2's conformal-time FLRW lift (pairs in all four directions)
is spanning. Its cosmic-time catalogue lift (spatial pairs only), and the same construction
for Kasner, has rank $9$: the lapse direction $g_{00}$ is unreachable. Adding one cancelling
time pair, $P_0=N_0=v$, makes it spanning without changing the metric.
\item \textbf{One-sector reductions.} For $\sigma=A\,\mathrm dt+B\,\mathrm dr$ the reduced equations
are exactly two projections of the Einstein tensor,
\begin{equation}\label{eq:reducedproj}
\mathcal E_A=-\sqrt{-g}\,(2AG^{tt}+2BG^{tr}),\qquad
\mathcal E_B=-\sqrt{-g}\,(2AG^{tr}+2BG^{rr}).
\end{equation}
Within static spherical symmetry they determine the solution (symmetric criticality): at their
solution $A^2=B^2=C/r$ every component of $G^{\mu\nu}$ vanishes, and with matter they reproduce
the TOV benchmarks. They do not govern perturbations outside the symmetry class.
\end{enumerate}
\end{proposition}
\noindent\textbf{Verification.} Item~1: Schwarzschild and Kerr ($a=0.9M$) at $r/M=10$, both horizons,
between them, and $r/M=0.05$; $\det\Sigma/v^4=1.00000000$, rank $10$, and metric reproduced to
round-off throughout. Item~2: rank $9$ with $g_{00}$ the unique unreachable direction for the
spatial lift; rank $10$ for the conformal-time lift ($a=10^{-3}$--$3$) and for the augmented lift
($a=10^{-3}$--$50$). Item~3: both ratios in Eq.~\eqref{eq:reducedproj} computed symbolically as exactly
$-1$, and all ten components of $G^{\mu\nu}$ vanish symbolically at the reduced solution.

\noindent\textbf{Consequence.} Every background used in Chapters~2--4 has a spanning QGD
representation, so by Theorem~\ref{thm:dof} perturbations about black holes and cosmologies
carry exactly the dynamics of general relativity. The one-sector and spatial-only
representations remain valid for the backgrounds themselves but must be replaced by their
spanning lifts for perturbation theory.

\section{What changes in the unrestricted derivative-dependent metric map}
\label{sec:derivativemap}
The master proposal~\eqref{eq:master} permits a nonzero term
$-\lambda\mathcal Q_{ab}$. Its functional form and coefficient are
currently unmatched; the existence of a symbol $\mathcal Q_{ab}$
does not establish the independent field equation $\square_g\sigma_\mu$
or any particular force correction. A conditional chain rule can,
however, be proved \emph{before} choosing that operator.

\begin{proposition}[Functional chain rule for a first-derivative metric map]
\label{prop:jet}
Hold the reference coframe fixed and suppose the composite metric is a
smooth local function $g_{\mu\nu}(\sigma,\partial\sigma)$ of the
independent natural-invariant fields and their first coordinate
derivatives. Assume the matter and gravitational action depend on
$\sigma$ only through the resulting metric, and its metric Euler
variation is
$\delta\Gamma=\int\dd^4x\,\sqrt{-g}\,\mathcal E_g^{\mu\nu}\,
\delta g_{\mu\nu}$ after the appropriate boundary treatment.
Then, for arbitrary compactly supported $\delta\sigma$, the
natural-invariant Euler derivative obeys
\begin{equation}\label{eq:jet}
\boxed{\displaystyle
\frac{\delta\Gamma}{\delta\sigma^{(s)}_\rho}
=\sqrt{-g}\,\mathcal E_g^{\mu\nu}
\frac{\partial g_{\mu\nu}}{\partial\sigma^{(s)}_\rho}
-\partial_\alpha\!\left(
\sqrt{-g}\,\mathcal E_g^{\mu\nu}
\frac{\partial g_{\mu\nu}}
{\partial(\partial_\alpha\sigma^{(s)}_\rho)}\right).}
\end{equation}
If $g$ depends on second or higher derivatives of $\sigma$, the
corresponding alternating higher total derivatives are required.
Direct nonmetric $\sigma$ couplings add their \emph{own} Euler
functional derivatives.
\end{proposition}
\begin{proof}
Use the ordinary chain rule on the specified local metric map,
substitute into the metric Euler variation, and integrate by parts
the term proportional to
$\partial_\alpha\delta\sigma^{(s)}_\rho$. Compact support removes
the boundary contribution. A second integration-by-parts argument
handles second derivatives, and so on.
\end{proof}
This equation expresses \emph{how} a well-defined derivative metric
operator would enter; it does not determine an independent kinetic
coefficient from the already accepted action. Without the actual
$\mathcal Q_{ab}$ and its admissible boundary terms, no numerical
quantum-force coefficient can be assigned to this contribution.

% === QGD-101 insertion, corrected QGD-105 ===
\begin{remark}[The derivative channel on Chapter~2's constructed branch (QGD--101/105)]
\label{rem:Q-absorbed}
For both trial completions of $\mathcal Q$ in Chapter~2 and any $\lambda\neq0$,
Chapter~2 constructs a particular nine-sector branch on which the metric map has
a differential right inverse (Theorems on the nine-sector lift). On that branch,
a metric-mediated action gives exactly the effective metric equations, whatever
$\lambda$ is. As Chapter~2 states, this concerns a particular spanning branch,
not all field configurations; on other branches, or with $\sigma$ dependence not
passing through $g$, the question remains open.
\end{remark}
% === end QGD-101/105 insertion ===

\section{Does this equation say anything beyond GR?}
\label{sec:equivalence}
One projected field equation per covector need not contain all ten
independent symmetric metric equations. What matters is the rank of
the map from independently variable signed covectors to the ten
metric components.

\begin{proposition}[Exact regular signed section (D)]
\label{prop:section}
In a fixed local reference frame let $h_{\mu\nu}
=g_{\mu\nu}-\bar g_{\mu\nu}$ be any real symmetric $4\times4$
matrix. For any constant $a>0$ set
\begin{equation}\label{eq:PN}
\boxed{P=aI_4+\frac{h}{4a},\qquad
N=aI_4-\frac{h}{4a}.}
\end{equation}
The four rows of $P$ and four rows of $N$ define eight signed
covectors, and
\begin{equation}\label{eq:identity}
\boxed{P^{\mathsf T}P-N^{\mathsf T}N=h.}
\end{equation}
For every symmetric variation $\delta h$, the choice
$\delta P=\delta h/(4a)$ and $\delta N=-\delta h/(4a)$
yields $\delta g=\delta h$. Thus the metric differential has rank ten
on this section.
\end{proposition}
\begin{proof}
Expand both products in \eqref{eq:identity}; the $a^2I$
and $h^2/(16a^2)$ terms cancel, leaving $h$.
Differentiation gives the displayed right inverse.
\end{proof}

\begin{theorem}[Conditional equivalence with Einsteinian dynamics (D$\mid$A)]
\label{thm:equiv}
Suppose the full effective action depends on the signed covectors
\emph{only through} $g(\sigma)$, and the admissible covector
variations possess a local right inverse to arbitrary symmetric
metric variations, as on the section \eqref{eq:PN}. Then
\begin{equation}\label{eq:equivalence}
\boxed{\cS_{(s)}^\rho=0\ \text{for every sector }s
\quad\Longleftrightarrow\quad
\cE_g^{\mu\nu}=0.}
\end{equation}
Thus, on this regular full-rank branch, the action-derived
natural-invariant equations are equivalent to the \emph{specified}
effective Einstein--matter equations.
\end{theorem}
\begin{proof}
If $\cE_g=0$, every projected Euler derivative is zero.
Conversely, let every covector Euler derivative vanish.
Use the differential right inverse of Proposition~\ref{prop:section}
to express any compactly supported symmetric metric test variation
as an admissible signed-covector variation. Then
$\delta\Gamma=0$ for all metric tests and the fundamental lemma
of the calculus of variations gives $\cE_g=0$.
\end{proof}

\noindent\textbf{Interpretation.}
Writing the Einsteinian equations using natural invariants is a
nontrivial local representation and an alternative choice of variables.
It does \emph{not}, by itself, generate additional observables beyond
GR. In the metric-only sector the transformation
$S\mapsto O(x)S$ with $S=\binom{P}{N}$ and
$O^{\mathsf T}\mathrm{diag}(I_4,-I_4)O=\mathrm{diag}(I_4,-I_4)$
leaves $g$ unchanged. The resulting local $O(4,4)$ freedom is a
parametrization redundancy wherever it is a symmetry of the complete
action and matter couplings. A physical $\sigma$ representative need
not be unique; adding direct interactions could change this symmetry.

\section{The calibrated gravitational field and the normalization distinction}
\label{sec:connection}
Chapter~3 independently varies both components $A(r)$ and $B(r)$ of
the declared classical spherical one-form $\sigma=A\,\dd x^0+B\,\dd r$,
without setting $A=B$ in the action. On its nondegenerate, asymptotically
flat vacuum branch, the adopted Einstein--Hilbert sector gives
\begin{equation}\label{eq:classical}
\boxed{\displaystyle A(r)^2=B(r)^2=\frac{C_g}{r},\qquad
 C_g=\frac{2G_NM}{c^2}=:r_s.}
\end{equation}
The value of $C_g$ is fixed \emph{once} by the Newtonian exterior
acceleration. It is not derived from the microscopic Dirac equation;
its calibration is legitimate correspondence with the observed
Newtonian mass parameter, as in the historical Einstein normalization.
No other coefficient is adjusted to recover higher orders.

When an independently selected reference fermion mass $m_{\rm ref}>0$
supplies a reduced Compton length $\bar\lambda_{\rm ref}=\hbar/(m_{\rm ref}c)$,
Chapter~3 defines the dimensionless \emph{geometric inverse-strength
convention}
\begin{equation}\label{eq:alphageom}
\boxed{\displaystyle |\alpha_G^{\rm geom}|^2:=\frac{\bar\lambda_{\rm ref}}{C_g}
=\frac{\hbar c}{2G_NMm_{\rm ref}}.}
\end{equation}
The original common spinor multiplier $\alpha_G^{\rm spin}$ in
$\Psi=(\alpha_G^{\rm spin}/r)\sum_a\psi_a u_a e^{\pm iS_a/\hbar}$
is \emph{not} thereby normalized from the geometric equations.
Indeed, the simultaneous change
$\alpha_G^{\rm spin}\mapsto z\alpha_G^{\rm spin}$,
$\psi_a\mapsto\psi_a/z$ with $z\ne0$ leaves the physical state
unchanged. Its absolute normalization comes from an independently
specified current or occupied-state convention. In particular,
$\alpha_G^{\rm spin}=\alpha_G^{\rm geom}$ is not a theorem of
Chapters~1--3.

The observable static Killing-energy potential per unit held test
mass is
\begin{equation}\label{eq:pot}
\boxed{\displaystyle\Phi_{\rm stat}(r)=c^2\left(\sqrt{1-\frac{C_g}{r}}-1\right),
\qquad r>C_g.}
\end{equation}
Its convergent, wholly \emph{classical} expansion is
\begin{equation}\label{eq:classseries}
\boxed{\displaystyle\Phi_{\rm stat}(r)=-\frac{G_NM}{r}
-\frac{G_N^2M^2}{2c^2r^2}
-\frac{G_N^3M^3}{2c^4r^3}
-\frac{5G_N^4M^4}{8c^6r^4}+\cdots .}
\end{equation}
All powers of $\hbar$ cancel when \eqref{eq:alphageom} is used
merely as a parametrization of $C_g$. Thus these higher orders are
not evidence of quantum-corrected vacuum gravity.

\section{The microscopic Dirac bridge: what it fixes}\label{sec:bridge}
For one locally on-shell, positive-energy, outward-propagating Dirac
WKB branch with conserved particle flux $\dot N>0$, Chapter~1 and
Chapter~3 give the magnitude-only current identity
\begin{equation}\label{eq:fluxbridge}
\boxed{\displaystyle |\alpha_G^{\rm spin}\psi_0(r)|^2
=\frac{\dot N\,\Delta(r)}{8\pi c^2p(r)},\qquad
\Delta=E_{\rm loc}+m_{\rm eff}c^2,\quad p>0.}
\end{equation}
This is a transport constraint. It determines the envelope only
relative to an independently fixed momentum profile and flux; it does
\emph{not} determine a $1/r$ gravitational potential, the geometric
normalization, or an absolute phase. A static bound source with
zero net radial current is \emph{not} normalized using this outgoing
flux formula. Its spinor profiles must be obtained from its own
coupled boundary-value equations with conserved occupation.

\section{A controlled quantum matter contribution to the source}\label{sec:quantumsource}
To make the $\hbar$ dependence precise without silently inserting a
new gravitational operator, use a stated positive-energy,
fixed-spin, nonrelativistic Dirac-to-Pauli limit and write
$\varphi=\sqrt n\,e^{iS/\hbar}\chi$ with constant unit Pauli
spinor $\chi$, $n>0$, and fermion mass $m_f>0$.
The polar identity for the Schr\"odinger kinetic energy is
\begin{equation}\label{eq:quantumenergy}
\boxed{\displaystyle E_{\rm kin}[n,S]
=\int\dd^3x\left[
\frac{n|\bm\nabla S|^2}{2m_f}
+\frac{\hbar^2}{8m_f}\frac{|\bm\nabla n|^2}{n}\right].}
\end{equation}
The quantum density-gradient functional and associated momentum-flux
tensor (one admissible improvement convention) satisfy
\begin{equation}\label{eq:qstress}
\boxed{\displaystyle Q[n]=-\frac{\hbar^2}{2m_f}
\frac{\nabla^2\sqrt n}{\sqrt n},\qquad
\Pi^Q_{ij}=\frac{\hbar^2}{4m_f}
\left(\frac{\partial_i n\partial_j n}{n}-\partial_i\partial_j n\right),\qquad
\partial_j\Pi_{ij}^Q=n\partial_i Q.}
\end{equation}
This is a derived coefficient of \emph{quantum matter stress},
not a derived coefficient of a standalone
$\hbar^2(\nabla\sigma)^2$ gravitational kinetic term.
To leading nonrelativistic order, the \emph{adopted} sourced
Einstein--Hilbert equations reduce to Poisson's law:
\begin{equation}\label{eq:Uquantum}
\boxed{\displaystyle
\nabla^2\Phi=4\pi G_Nm_fn,\qquad
\mathcal U:=-\frac{2\Phi}{c^2},\qquad
\nabla^2\mathcal U=-\frac{8\pi G_Nm_f}{c^2}n.}
\end{equation}
If the physical source profile $n(x;\hbar)$ is independently determined
from its quantum dynamics, its $\hbar$ dependence influences the
sourced metric and consequently the corresponding temporal
natural-invariant \emph{combination}. This is a source effect, not a
new vacuum propagation law. Relativistic stress, pressure, spin,
multi-orbital occupations and binding require the complete
Dirac--Einstein source instead of simply substituting one $n$.

\begin{proposition}[Compact spherical exterior at fixed mass (D$\mid$A)]
\label{prop:exterior}
Let $\rho_\hbar(r)$ be an integrable, nonnegative spherically symmetric
mass density with fixed total mass
$M=4\pi\int_0^\infty\rho_\hbar(r)r^2\,\dd r$, compactly supported
in $r\leq R$. Assume Poisson's equation and $\Phi(\infty)=0$.
Then $\Phi(r)=-G_NM/r$ and
$F_r=-G_NMm_t/r^2$ for every $r>R$, independently of the interior
$\hbar$ dependence.
\end{proposition}
\begin{proof}
Spherical Poisson integration gives
$r^2\Phi'(r)=4\pi G_N\int_0^r\rho_\hbar(s)s^2\,\dd s$.
The integral equals $M/(4\pi)$ outside the support. Integrate once
more with the stated asymptotic boundary condition.
\end{proof}
Thus a universal exterior term proportional to
$G_NM\hbar^2/(m_f^2c^2r^4)$ does \emph{not} follow from changing only
the profile of a compact spherical quantum source while keeping its
mass and the Poisson law fixed. Noncompact sources, altered total mass,
non-spherical multipoles or independently derived quantum field
equations are distinct questions.
% === QGD-113 insertion ===
\begin{remark}[Varying spin (QGD--113)]
The fixed-spin assumption above is a special case. When the spin direction
$\hat s$ varies, the exact spin-$\frac12$ flow law of Section~\ref{sec:spinhalf} adds the
spin quantum potential $(\hbar^2/8m_f)|\nabla\hat s|^2$ and a spin-vorticity force; both
vanish for constant $\hat s$, recovering Eq.~\eqref{eq:qstress}.
\end{remark}
% === end QGD-113 insertion ===

\section{Status of the historical quantum-force candidate}\label{sec:historicalforce}
The proposed universal attractive force
\begin{equation}\label{eq:historic9}
\boxed{\displaystyle F_{r,\rm hist}=-\frac{G_NMm_t}{r^2}
\left[1+\frac92\left(\frac{\hbar}{m_fc r}\right)^2+\cdots\right]}
\end{equation}
remains a historical \emph{candidate}, not a conclusion of the
accepted action. The reciprocal of $e^{2im_fcr/\hbar}$ is
$e^{-2im_fcr/\hbar}$: inverting nonconstant Taylor terms separately
does not invert the complete exponential. The historical cubic energy
also gives a sign opposite to an enhancement of attraction.

A first-order weak-field expansion of a specified Liouville-normal,
separated Schwarzschild radial Dirac equation produces coefficients
depending on its angular channel $\kappa$. Their derivatives are
not coordinate-, tetrad-, or field-redefinition-invariant force
observables. In particular the lowest $\kappa=-1$ channel does not
produce the historical term at order $G_N\hbar^2/r^4$ in that normal
form. Nothing here establishes a universal quantum change to the
exterior gravitational acceleration. To claim one, first derive a
covariant $\hbar$-dependent action, vary it including all connection
and metric dependence, and calculate a gauge-/representation-invariant
physical observable.



\subsection{Independent referee test of the radial Dirac coefficient (QGD--098)}
\label{sec:radial-referee}
The following is a \emph{conditional theorem about a quantum test field}.
It does not modify the natural-invariant field equations. Adopt the
classical, Newton-calibrated Schwarzschild exterior of
Eq.~\eqref{eq:classical}, minimally couple a massive four-component
Dirac field through the induced orthonormal-frame spin connection,
and use the standard separated, fixed-tetrad radial basis. This
additional coupling and the reduction to a single radial mode are
assumptions of the \emph{test-field model}, not consequences of the
flat-space spinor-component ratios alone. Keep $r>r_s>0$ and define
\begin{equation}\label{eq:radial-parameters}
\ell_C:=\frac{\hbar}{m_fc},\quad x:=\frac r{\ell_C},\quad
 a:=\frac{r_s}{\ell_C},\quad e:=\frac{E}{m_fc^2},\quad
 F:=1-\frac ax,\quad \frac{\dd}{\dd x_*}:=F\frac{\dd}{\dd x},
\end{equation}
where $\kappa\in\mathbb Z\setminus\{0\}$ is the spinor angular
index. For the stated radial spinor gauge the first-order field
equations read
\begin{equation}\label{eq:radial-first-order}
\boxed{\begin{aligned}
\frac{\dd f}{\dd x_*}&=(e+\sqrt F)g-\frac{\kappa\sqrt F}{x}f,\\
\frac{\dd g}{\dd x_*}&=\frac{\kappa\sqrt F}{x}g+(\sqrt F-e)f.
\end{aligned}}
\end{equation}
\begin{proposition}[Liouville-normal radial coefficient (D$\mid$A)]
\label{prop:radial-test}
Assume Eqs.~\eqref{eq:radial-parameters}--\eqref{eq:radial-first-order},
$e=1$, and $a/x<1$. Put $B:=1+\sqrt F$, $K:=\kappa\sqrt F/x$, and
\begin{equation}\label{eq:radial-pq}
P:=\frac{F_x}{F}-\frac{B_x}{B},\qquad
Q:=F K_x-F\frac{B_x}{B}K-K^2+1-F.
\end{equation}
Then elimination of $g$ gives
\begin{equation}\label{eq:radial-second-order}
f_{xx}+P f_x+\frac Q{F^2}f=0.
\end{equation}
After the invertible local Liouville rescaling
$f=\exp(-\tfrac12\int P\,\dd x)u$, the transformed equation is
\begin{equation}\label{eq:radial-liouville}
\boxed{u_{xx}+\mathcal V_\kappa(x,a)u=0,\qquad
\mathcal V_\kappa=\frac Q{F^2}-\frac12P_x-\frac14P^2.}
\end{equation}
Its expansion at fixed $x>0$ about $a=0$ is
\begin{equation}\label{eq:radial-series-general}
\boxed{\begin{aligned}
\mathcal V_\kappa={}&-\frac{\kappa(\kappa+1)}{x^2}
+\frac ax+\frac{3-4\kappa^2-\kappa}{4}\frac a{x^3}\\
&+\frac{2a^2}{x^2}
+\frac{-64\kappa^2+4\kappa+69}{64}\frac{a^2}{x^4}
+O(a^3).
\end{aligned}}
\end{equation}
In particular,
\begin{equation}\label{eq:radial-kminus1}
\boxed{\mathcal V_{-1}(x,a)=\frac ax+\frac{2a^2}{x^2}
+\frac{a^2}{64x^4}+O(a^3),}
\end{equation}
so the $O(a/x^3)$ correction vanishes in the $\kappa=-1$
channel. The displayed expansions are local analytic expansions
under the stated hypotheses; they are not uniform near $r=r_s$.
\end{proposition}
\begin{proof}
Solve the first equation in \eqref{eq:radial-first-order} for
$g=(Ff_x+Kf)/B$. Differentiate with respect to $x_*$, substitute the
second equation, and collect the $f_x$ and $f$ coefficients to
obtain Eq.~\eqref{eq:radial-second-order}. Direct differentiation
of the Liouville rescaling gives Eq.~\eqref{eq:radial-liouville}.
Using $\sqrt{1-a/x}=1-a/(2x)-a^2/(8x^2)+O(a^3)$ in the resulting
\emph{exact} coefficient gives Eq.~\eqref{eq:radial-series-general}.
Setting $\kappa=-1$ proves Eq.~\eqref{eq:radial-kminus1}.
\end{proof}
\begin{remark}[What is and is not a quantum force]
\label{rem:radial-not-force}
At $e=1$ one may \emph{formally} rewrite
Eq.~\eqref{eq:radial-liouville} as a zero-energy one-dimensional
Schr\"odinger equation, with
$U_{\rm radial}:=-m_fc^2\mathcal V_\kappa/2$.
Its $O(a/x^3)$ contribution is
\begin{equation}\label{eq:radial-gamma}
\delta U_\kappa=-\gamma_\kappa
\frac{G_NM\hbar^2}{m_fc^2r^3},\qquad
\gamma_\kappa:=\frac{3-4\kappa^2-\kappa}{4}.
\end{equation}
Differentiating this \emph{representation-dependent radial term} yields
$-3\gamma_\kappa G_NM\hbar^2/(m_fc^2r^4)$, but this is
\textbf{not} by itself a local force observable or an invariant
correction to the gravitational metric. In particular,
$\gamma_{-1}=0$ and $\gamma_{+1}=-1/2$. The historical 
\emph{attractive} coefficient $9/2$ would require
$\gamma_\kappa=3/2$, i.e.
$4\kappa^2+\kappa+3=0$, which has no real root.
The $a^2/(64x^4)$ term at $\kappa=-1$ is likewise an
\emph{effective radial-equation coefficient}, not a new vacuum
Einstein tensor. It must not be reported as a distinct QGD
quantum-gravitational force.
\end{remark}
\paragraph{Referee/proof boundary.}
The result above rests on a stated classical Schwarzschild exterior,
the minimally coupled test Dirac equation, a fixed tetrad and spinor
angular basis, and a specified threshold energy $e=1$. It therefore
follows from Chapters~1--2 only \emph{after} their declared effective
metric/Dirac coupling is adopted. It is not derived solely from the
flat-space macroscopic spinor ratios. Exact symbolic differentiation
verifies Eqs.~\eqref{eq:radial-second-order}--\eqref{eq:radial-kminus1};
this is not a Lean-kernel proof. A physical quantum-gravity claim
would additionally require the modified, conserved metric/source
field equations and an invariant observable, with no further fit
to the historical $9/2$ coefficient.


\section{A published semiclassical Dirac benchmark in natural-invariant variables}
\label{sec:diracbenchmark}
The historical $9/2$ Compton-scale force is not established by the
original phase manipulation.  There is, however, a distinct and
quantitatively fixed $\hbar$-dependent \emph{gravitational} correction
that can be imported into the effective metric sector.  Carlson,
Anderson, Fabbri, Fagnocchi, Hirsch and Klyap computed the leading
far-field renormalized stress tensor of one free \emph{massless
four-component} Dirac field in an asymptotically empty (Boulware)
state, and checked its gravitational backreaction against previous
one-loop calculations~\cite{Carlson2010}.  This is a benchmark of
semiclassical Einstein--Dirac physics, not an independent derivation
of gravitational dynamics from the original flat-space QGD Dirac action.

\begin{proposition}[Massless Dirac far-field gravitational correction
(D$\mid$A)]\label{prop:diracfarfield}
Adopt the metric-only effective Einstein--Hilbert action with a
minimally coupled quantized free massless four-component Dirac field,
the renormalized Boulware state, asymptotic flatness, fixed ADM mass
$M>0$, areal radius $r$, and leading-order semiclassical backreaction.
In units $G_N=c=\hbar=1$, the published leading stress tensor is
\begin{equation}\label{eq:diracRSET}
\boxed{\left\langle T^\mu{}_{\nu}\right\rangle_{\rm ren}
=\frac{M}{40\pi^2r^5}\,
\operatorname{diag}(4,2,-3,-3)+O(r^{-6}).}
\end{equation}
For the far-field static metric
\begin{equation}\label{eq:scmetric}
\dd s^2=-f(r)\dd t^2+
\left(1-\frac{2\mathfrak m(r)}r\right)^{-1}\dd r^2
+r^2\dd\Omega^2,\qquad
f(r)=\left(1-\frac{2\mathfrak m(r)}r\right)e^{2\nu(r)},
\end{equation}
the asymptotically normalized functions satisfy
\begin{align}
\mathfrak m(r)&=M\left(1+\frac1{5\pi r^2}\right)
+O(r^{-3}),\label{eq:massfn}\\
\nu(r)&=\frac{M}{15\pi r^3}+O(r^{-4}),\label{eq:nufn}\\
f(r)&=1-\frac{2M}{r}-\frac{4M}{15\pi r^3}
+O(r^{-4}),\label{eq:scf}
\end{align}
where in the last expression $O(r^{-4})$ also includes mixed
classical--semiclassical orders.  Restoring $G_N,c,\hbar$, the
\emph{leading far-field} supported-static-body Killing-energy
potential per unit test mass is
\begin{equation}\label{eq:QGD97Phi}
\boxed{\begin{aligned}
\Phi_{\rm QGD}^{\rm sc}(r)={}&c^2\left(
\sqrt{1-\frac{2G_NM}{c^2r}}-1\right)\\
&-\frac{2G_N^2M\hbar}{15\pi c^3r^3}
+O\!\left(\frac{G_N^3M^2\hbar}{c^5r^4}\right)
+O(\hbar^2),
\end{aligned}}
\end{equation}
where the final $O(\hbar^2)$ denotes omitted higher quantum orders,
not a uniform SI-dimensional bound.  The displayed quantum coefficient
is meaningful in the far field $r\gg r_s=2G_NM/c^2$, where the
quantum backreaction is perturbatively small.  It must not be
extrapolated unchanged to $r\simeq r_s$.
\end{proposition}
\begin{proof}
For the ansatz \eqref{eq:scmetric}, the radial Einstein equations
at leading quantum order read
\begin{equation}\label{eq:scEFE}
\mathfrak m'=-4\pi r^2\langle T^t{}_t\rangle_{\rm ren},
\qquad
\nu'=\frac{4\pi r
(\langle T^r{}_r\rangle_{\rm ren}
-\langle T^t{}_t\rangle_{\rm ren})}
{1-2\mathfrak m/r}.
\end{equation}
Substituting \eqref{eq:diracRSET} gives
$\mathfrak m'=-2M/(5\pi r^3)$ and
$\nu'=-M/(5\pi r^4)+O(r^{-5})$.
Integrating with $\mathfrak m(\infty)=M$ and $\nu(\infty)=0$
proves \eqref{eq:massfn}--\eqref{eq:nufn}.  Consequently
$f=1-2M/r-4M/(15\pi r^3)+O(r^{-4})$.
For an observer held static in this exterior,
$\Phi=c^2(\sqrt{f}-1)$; expanding to first order in the
renormalized Dirac stress and restoring dimensions proves
\eqref{eq:QGD97Phi}.  This is a sourced asymptotic expansion,
not the solution of a freely chosen scalar force law.
\end{proof}

\begin{remark}[Quantum gravitational force and Planck-scale residue]
The \emph{quantum part} of the signed radial force on a held static
test mass $m_t$ has leading far-field behavior
\begin{equation}\label{eq:scforce}
\boxed{\delta F_r^{\rm sc}
=-\frac{2G_N^2Mm_t\hbar}{5\pi c^3r^4}
+\text{higher orders}.}
\end{equation}
Relative to the Newtonian attraction, its leading magnitude is
$\frac{2}{5\pi}\frac{G_N\hbar}{c^3r^2}$,
which is governed by the Planck length squared, not the inverse
square of an arbitrarily selected probe mass.  It is therefore
\emph{not} the historical $(9/2)(\hbar/(m_fc r))^2$ correction.
The classical piece $-m_t\Phi_{\rm stat}'$ remains defined by the
exact background square root.  The quantum term is a controlled
far-field correction, not an exact all-radius force.
\end{remark}

\begin{remark}[Natural-invariant reconstruction and status]
The full signed/tetrad QGD master metric \eqref{eq:master} is
unchanged.  For the nondegenerate one-sector static spherical
\emph{metric representation} discussed in Chapter~2, write
$\sigma=A(r)\dd t+B(r)\dd r$ and put
$F=1-A^2$, $D=1+B^2-A^2$ so that the diagonal radial metric
coefficient is $D/F$.  Reconstruct the semiclassical metric by
$A^2=1-f$ and $D/F=(1-2\mathfrak m/r)^{-1}$.
Then, in the far field and to the first displayed quantum order,
\begin{equation}\label{eq:scsigmaAB}
\boxed{\begin{aligned}
A^2(r)&=\frac{2M}{r}+\frac{4M}{15\pi r^3}
+O(r^{-4}),\\
B^2(r)&=\frac{2M}{r}+\frac{2M}{5\pi r^3}
+O(r^{-4})
\end{aligned}}
\qquad(G_N=c=\hbar=1).
\end{equation}
Thus the familiar equality $A^2=B^2$ is a property of the
classical Schwarzschild branch, \emph{not} of the quantum-corrected
source.  This construction shows that QGD variables can
represent the published semiclassical metric; it does not
establish a new dynamical law for $\sigma_\mu$.
\end{remark}

\paragraph{Assumptions and formal-verification boundary.}
The algebraic implications from \eqref{eq:diracRSET} and
\eqref{eq:scEFE} to \eqref{eq:QGD97Phi} are suitable for a
Lean-style theorem once analytic asymptotic bounds and positivity
conditions on the metric are specified.  The renormalized-stress
coefficient in \eqref{eq:diracRSET} is a \emph{published quantum
field-theoretic input}, not a theorem derived from Chapter~1.
The complete original four-mode wavefunction has not been proved
to generate this Boulware vacuum as its unique quantum state.
At the current metric-only effective level, the equivalence theorem
in the next section applies: this correction is already predicted
by semiclassical GR with the same Dirac matter.


\section{Can QGD predict a quantum force independently of semiclassical GR?}
\label{sec:independent}
This is the central discriminating question of QGD--096. Here
\emph{semiclassical Einstein gravity} means a classical metric sourced by
a renormalized expectation value or effective action of quantum matter;
\emph{quantized gravitational effective field theory} includes
fluctuations of the gravitational field itself. The two comparisons are
different. A Dirac wave in a prescribed metric, the quantum-potential
term of a Pauli reduction, and an $\hbar$-dependent source profile are
not, by themselves, quantum corrections to the gravitational vacuum
response. The relevant test is an invariant observable calculated from
an effective dynamics which QGD derives but which the selected GR-plus-
quantum-matter comparator does not already possess.

\begin{theorem}[Metric-only quantum-equivalence obstruction (D$\mid$A)]
\label{thm:quantumno-go}
Let $\mathcal F$ be the algebraic signed-sector map $\sigma\mapsto
g[\sigma]$ on a regular branch whose differential is surjective onto
symmetric metric variations, with identical boundary conditions and
matter states on both sides. Let $\Gamma_{\rm sc}[g,\Psi]$ be a
specified semiclassical Einstein--matter effective functional,
possibly including renormalized quantum matter loops, and assume
\begin{equation}\label{eq:qgd-sc-pullback}
\boxed{\Gamma_{\rm QGD}[\sigma,\Psi]
=\Gamma_{\rm sc}[\mathcal F(\sigma),\Psi]}
\end{equation}
with no independent $\sigma$ interaction or separately quantized
$\sigma$ measure. Then stationary QGD fields on this branch and
stationary fields of $\Gamma_{\rm sc}$ determine the same set of
physical metrics and metric/matter observables. In particular, no
additional universal $\hbar$-dependent gravitational-force coefficient
is implied solely by expressing $g$ through $\sigma$.
\end{theorem}
\begin{proof}
For any admissible variation, the functional chain rule gives
\begin{equation}\label{eq:quantum-pullback-chain}
\delta\Gamma_{\rm QGD}
=\int\dd^4x\,\sqrt{-g}\,\mathcal E_{\rm sc}^{\mu\nu}
\frac{\partial g_{\mu\nu}}{\partial\sigma^{(s)}_\rho}
\delta\sigma^{(s)}_\rho
+\delta_\Psi\Gamma_{\rm sc},
\end{equation}
where repeated sector indices are summed. A stationary metric of
$\Gamma_{\rm sc}$ therefore makes the QGD pullback stationary.
Conversely, by surjectivity every compactly supported symmetric
$\delta g$ is induced by an admissible $\delta\sigma$; stationarity
of the pullback and the independently imposed matter equations imply
$\mathcal E_{\rm sc}^{\mu\nu}=0$. Because all specified observables
are functionals of the same physical metric and matter state, they
agree. The result is an on-shell, mean-effective-action theorem: it
makes no claim about quantizing the redundant variables with a new
functional measure, about rank-deficient branches, or about a
new action not covered by \eqref{eq:qgd-sc-pullback}.
\end{proof}

\begin{remark}[What the theorem does, and does not, rule out]
It rules out \emph{new observable coefficients arising from a metric
reparametrization alone}. It does not forbid a distinct QGD quantum
mechanism. For example, an independently matched covariant correction
$\Delta\Gamma_{\rm QGD}[\sigma,\Psi;\hbar]$ could alter the field
equations, whether through additional natural-invariant physics, a
quantum-derived metric operator, or a consistently specified quantum
measure. But a coefficient copied from a known GR effective action,
or chosen to reproduce the historical force, would not constitute an
independent QGD prediction. A meaningful comparison must specify
the same physical source, asymptotic conditions, renormalization
prescription, and gauge-invariant observables.
\end{remark}

\subsection{An exact exterior test of the historical $9/2$ law}
\label{sec:poisson-residual}
The attractive force candidate in Eq.~\eqref{eq:historic9} can be
integrated unambiguously at the displayed order. Put
$\ell_C=\hbar/(m_fc)$, where the reference species $m_f$ must be
specified independently. In the weak-field nonrelativistic
interpretation, $F_r=-m_t\partial_r\Phi$ and $\Phi(\infty)=0$ imply
\begin{equation}\label{eq:quantum-candidate-potential}
\boxed{\displaystyle
\Phi_{9/2}(r)=-\frac{G_NM}{r}
-\frac{3G_NM\ell_C^2}{2r^3}
+\text{higher-order terms}.}
\end{equation}
For the isolated displayed $r^{-3}$ term and $r>0$,
$\nabla^2(r^{-3})=6r^{-5}$, so
\begin{equation}\label{eq:quantum-laplacian-test}
\boxed{\displaystyle
\nabla^2\Phi_{9/2}(r)=-\frac{9G_NM\ell_C^2}{r^5}\ne0
\quad (r>R),}
\end{equation}
where $R$ is the support radius of an assumed compact source and
higher-order terms have been omitted consistently. If one tried to
retain the \emph{unmodified} Poisson equation, the same radial tail
would demand an exterior effective density
\begin{equation}\label{eq:quantum-rho-eff}
\boxed{\displaystyle
\rho_{\rm eff}(r)=-\frac{9M\ell_C^2}{4\pi r^5},
\qquad r>R,}
\end{equation}
which is not the original empty exterior. This is a diagnostic
\emph{effective} density, not a claim about negative physical matter.
Therefore the historical universal correction cannot be generated
solely by changing a compact spherical source profile under the same
vacuum Poisson equation; it requires a justified exterior quantum
source/response, a modified field equation, or other physics.

At the fully relativistic level, this weak-field calculation is not
by itself a complete covariant metric solution. The independently
varied Einstein--Hilbert vacuum branch of Chapter~2 already fixes
the exact spherical exterior to Schwarzschild. A new universal
$1/r^3$ metric perturbation must consequently be checked against the
\emph{complete} modified field equations and their conservation and
boundary conditions; a scalar force ansatz cannot perform that check.
If $m_f$ denotes the test particle mass, its Compton scale also makes
the force state/species-dependent, so the claim is not automatically a
universal exterior metric correction.

\subsection{A falsifiable QGD-specific matching test}
\label{sec:delta-action}
A future new term should be expressed as the difference between a
microscopically justified QGD effective functional and the specified
semiclassical comparator:
\begin{equation}\label{eq:delta-gamma}
\boxed{\displaystyle
\Delta\Gamma[\sigma,\Psi;\hbar]
:=\Gamma_{\rm QGD}[\sigma,\Psi;\hbar]
-\Gamma_{\rm sc}[g[\sigma],\Psi;\hbar].}
\end{equation}
At fixed physical matter state and boundary data, a nonzero
$\Delta\Gamma$ is \emph{necessary but not sufficient} for a new
observable: it may contain a boundary term, a redundancy, or an
operator removable by permissible redefinitions. The minimum
verification gates are:
\begin{enumerate}
\item Derive $\Delta\Gamma$ from the stated microscopic/macroscopic
 Dirac construction or declare any additional physical axiom.
\item Specify its covariant dependence on the full signed master
 metric, $\sigma^{(s)}$, and any derivative operators; include its
 independent normalization, units and boundary terms.
\item Vary the \emph{complete} functional, including connection,
 coframe, stress-energy and possible nonlocal contributions, and
 check Bianchi/conservation identities and regular branches.
\item Fix at most the permitted Newtonian calibration coefficient,
 then predict a second observable without re-fitting it.
\item Compare with \emph{both} semiclassical GR plus quantum matter
 and, where applicable, low-energy quantized-GR effective-field-theory
 corrections, at matched sources and renormalization convention.
\end{enumerate}
A quantum Dirac radial equation in the unmodified Schwarzschild metric
meets none of the first three gates for a new \emph{gravitational
field} operator. The previously derived quantum density-gradient
stress meets the microscopic matter-side test, but its gravitational
consequences within Eq.~\eqref{eq:action} are also consequences of
the chosen semiclassical Einstein--matter system. Hence the present
QGD action does \emph{not yet} close the deeper foundational gap of
demonstrating that QGD itself predicts a quantum correction
independently of semiclassical general relativity.


% ===================== QGD-110 INSERTION =====================
\section{The quantum structure of the master metric (QGD--110)}\label{sec:quantumstructure}
This section uses only the founding principles: F1, the microscopic Dirac field;
F2, coarse-graining and the macroscopic spinor equation; F3, the master metric
with its reference term; F4, the adopted Einstein--Hilbert-plus-matter action;
F5, the Newtonian calibration. Imported results are marked L.

\subsection{A metric-level quantum term and the field equations}
\begin{proposition}[Exact split identity (D)]\label{prop:split}
Write the physical metric as $g=g_{\rm cl}+q$. If $G[g]=8\pi G\,T$, then
identically
\begin{equation}\label{eq:qsplit}
G[g_{\rm cl}]=8\pi G\,(T+T_Q),\qquad T_Q:=-\frac{G[g]-G[g_{\rm cl}]}{8\pi G}.
\end{equation}
\end{proposition}
\noindent A quantum term in the master metric therefore always produces a
quantum-corrected Einstein equation for the classical part. By itself this is a
rewriting. It acquires content only when F1--F5 fix $q$, which they do through the
quantum stress of the matter.

\subsection{The quantum term in the static natural-invariant variables}
\begin{proposition}[Exact identity (D)]\label{prop:ABFH}
In the static branch ($F=1-A^2$, $D=1+B^2-A^2$, areal radius),
\begin{equation}\label{eq:ABFH}
A^2-B^2=1-FH,\qquad \frac{\mathrm d}{\mathrm dr}\ln(FH)=\frac{8\pi r(\rho+p_r)}{1-2\mathfrak m/r},
\end{equation}
with $H=D/F$. Hence $A=B$ exactly whenever $\rho+p_r=0$: in vacuum, for Maxwell
fields and for a cosmological constant.
\end{proposition}
\noindent For the vacuum of the microscopic Dirac field (Proposition~\ref{prop:diracfarfield}),
Eq.~\eqref{eq:scsigmaAB} gives
\begin{equation}\label{eq:ABquantum}
A^2-B^2=-\frac{2}{15\pi}\,\frac{G_N^2M\hbar}{c^5r^3},
\end{equation}
because vacuum polarization of the microscopic Dirac field has $\rho+p_r<0$. Inserting this $q$ into
Eq.~\eqref{eq:qsplit} returns $T_Q=-T^{\rm Dirac}$ in all three independent
components (verified symbolically). This is the quantum term F1 and F4 place in the
master metric. Its scale is $\ell_P^2$, not the historical $\ell_Q$, and by
Theorem~\ref{thm:quantumno-go} it coincides with semiclassical GR.

\begin{proposition}[Test of the historical derivative trial (D)]\label{prop:trialtest}
For the classical exterior sector $A=B=\sqrt{2M/r}$ on the flat spherical
reference, the trial $\mathcal Q_{ab}=(\bar\nabla_a\sigma^c)(\bar\nabla_b\sigma_c)$ has
only angular components, $\mathcal Q_{\theta\theta}=2M/r$, and in areal gauge gives,
to first order in $M$, $\Delta f=0$ and $\Delta h=-4\lambda M/R^3$. The derived term
requires $\Delta f=-4M\ell_P^2/(15\pi R^3)$ and $\Delta h=+2M\ell_P^2/(5\pi R^3)$.
No $\lambda$ matches both, so the trial is not the quantum term of the foundations.
\end{proposition}

\subsection{A quantum conformal factor derived from the wave equation}
\begin{proposition}[Quantum conformal factor (D)]\label{prop:qcf}
For a scalar wave $\psi=Re^{iS/\hbar}$ of mass $M$ obeying
$(\square_g-M^2c^2/\hbar^2)\psi=0$, the real part of the wave equation is exactly the
Hamilton--Jacobi equation of a conformally rescaled metric,
\begin{equation}\label{eq:qcf}
\tilde g^{\mu\nu}\partial_\mu S\partial_\nu S+M^2c^2=0,\qquad
\tilde g_{\mu\nu}=\Omega_\psi^2g_{\mu\nu},\qquad
\Omega_\psi^2=1-\lambda_C^2\frac{\square_gR}{R},\quad\lambda_C=\frac\hbar{Mc}.
\end{equation}
Wherever $\Omega_\psi^2>0$, the integral curves of $\tilde g^{\mu\nu}\partial_\nu S$
are timelike geodesics of $\tilde g$.
\end{proposition}
\begin{proof}
The real part is $g^{\mu\nu}\partial S\partial S+M^2c^2=\hbar^2\square_gR/R$ (verified
symbolically). For the second statement, $u_\mu=\partial_\mu S/(Mc)$ is a gradient
with $\tilde g(u,u)=-1$, so $u^\nu\tilde\nabla_\nu u_\mu=\frac12\partial_\mu\tilde g(u,u)=0$.
\end{proof}
\noindent This is the quantum potential of Chapter~1 (Eq.~wkb-quantumHJ) and of
Section~\ref{sec:quantumsource}, in geometric form. Calling the integral curves
\emph{trajectories} is an interpretation. The theorem itself is a statement about
the wave's phase and amplitude.

\subsection{How the quantum term enters the master metric}
\begin{theorem}[Two consistent routes, one inconsistent (D$\mid$A)]\label{thm:routes}
Under F1--F5:
\begin{enumerate}
\item \textbf{Per wave (consistent).} $\Omega_\psi^2$ multiplies the universal
master metric for that wave (spin-0 matter; spin-$\frac12$ matter in
Section~\ref{sec:spinhalf}),
$\tilde g^{(\psi)}_{\mu\nu}=\Omega_\psi^2T^a{}_\mu T^b{}_\nu[\eta_{ab}+\sum_ss_s\sigma^{(s)}_a\sigma^{(s)}_b]$,
and governs that wave's phase dynamics. Universal free fall is recovered as
$\lambda_C\to0$; at finite $\lambda_C$ the mass dependence is that of ordinary
quantum mechanics in a gravitational field.
\item \textbf{Through the source's stress tensor (consistent, universal).} On
shell the wave's stress tensor splits exactly as
\begin{equation}\label{eq:TQ}
T_{\mu\nu}=\rho\,u_\mu u_\nu+T^Q_{\mu\nu},\qquad
T^Q_{\mu\nu}=\hbar^2\bigl[\partial_\mu R\partial_\nu R-\tfrac12g_{\mu\nu}\bigl((\partial R)^2+R\,\square_gR\bigr)\bigr],
\end{equation}
with $\rho=M^2c^2R^2$ and the on-shell $u_\mu=\partial_\mu S/(Mc)$. Under F4, $T^Q$
curves spacetime with coupling $8\pi G$, and through Eq.~\eqref{eq:qsplit} it
appears as a universal quantum sector $q$ of the master metric. For a coherent
static source of size $L$, $T^Q_{00}/\rho c^2=O(\lambda_C^2/L^2)$.
\item \textbf{Source's conformal factor in the universal metric (inconsistent).}
For $g=(1+\varepsilon)g_{\rm cl}$, the linearized Einstein tensor
$G_{00}=-\nabla^2\varepsilon$ carries no factor of $G$; for a Gaussian amplitude it
equals $6\lambda_C^2/L^4$. This contradicts F4 at order $\hbar^2G^0$, and would
require replacing F4.
\end{enumerate}
\end{theorem}

\noindent The result is a two-level master metric derived from the founding
principles:
\begin{equation}\label{eq:twolevel}
\boxed{g_{\mu\nu}=T^a{}_\mu T^b{}_\nu\Bigl[\eta_{ab}+\sum_ss_s\sigma^{(s)}_a\sigma^{(s)}_b+q_{ab}\Bigr],
\qquad \tilde g^{(\psi)}_{\mu\nu}=\Omega_\psi^2\,g_{\mu\nu},}
\end{equation}
with sector signs $s_s=-\varepsilon_s$ in the notation of Eq.~\eqref{eq:master}, and
where $q$ is fixed by F4 from the quantum stress of the source (vacuum
polarization, Eq.~\eqref{eq:ABquantum}, and the coherent-wave stress $T^Q$), and
$\Omega_\psi$ by Eq.~\eqref{eq:qcf}. This gives the founding master metric's quantum
slot, historically $-\lambda\mathcal Q$, a derived occupant. As $\hbar\to0$, $q$'s
quantum part and $\Omega_\psi-1$ vanish, recovering the classical restriction and
every classical result of Chapters~2--4. The structure reproduces quantum
mechanics in curved spacetime and semiclassical gravity, and it does not quantize
the metric: the limitations of semiclassical gravity, for example for macroscopic
superpositions of a source, apply unchanged.

\begin{proposition}[No natural-invariant dynamics beyond $g$ under F1--F5 (D)]\label{prop:nobeyond}
The microscopic Dirac action uses a vierbein of $g$ and is locally Lorentz invariant (Chapter~2,
Clifford compatibility), so it depends only on $g$. The Maxwell, scalar-contact and
Einstein--Hilbert terms and the quantum matter effective action $\Gamma[g]$ are
functionals of $g$. Theorem~\ref{thm:quantumno-go} therefore applies: dynamics
acting on more than $g$ requires an additional founding principle.
\end{proposition}

\subsection{The per-wave law for spin-$\frac12$ matter (QGD--112, QGD--113)}\label{sec:spinhalf}
For spin-$\frac12$ matter no metric conformal to $g$ makes the flow lines geodesic
(QGD--111). The exact replacement follows from the macroscopic spinor equation
alone. Define $\sigma=\bar\psi\psi\neq0$, the convective momentum
$P_\mu=-\operatorname{Im}(\bar\psi\nabla_\mu\psi)/\sigma$, the spin remainder
$\Lambda=\nabla\bar\psi\cdot\nabla\psi-|\bar\psi\nabla\psi|^2/\sigma$, the spin tensor
$S^{ab}=\bar\psi\Sigma^{ab}\psi$ with $\Sigma^{ab}=\frac i2[\gamma^a,\gamma^b]$, and the spin
vorticity $F=\mathrm dP$.
\begin{theorem}[Spin-$\tfrac12$ flow law in a curved universal metric (D)]\label{thm:spinhalf}
In Chapter~1's conventions, $(\gamma^\mu\nabla_\mu)^2=-\square+\frac14R$ and
$[\nabla_\mu,\nabla_\nu]=-\frac14R_{ab\mu\nu}\gamma^a\gamma^b$ on spinors. Consequently, exactly,
\begin{equation}\label{eq:spinhalfflow}
P\cdot P+M^2+\tfrac14R=\frac{\square\sqrt\sigma}{\sqrt\sigma}-\frac\Lambda\sigma,\qquad
F_{\mu\nu}=-\frac{\operatorname{Im}(\nabla_{[\mu}\bar\psi\nabla_{\nu]}\psi)}{\sigma}
-\frac{\partial_{[\mu}\sigma\,P_{\nu]}}{\sigma}-\frac{R_{ab\mu\nu}S^{ab}}{4\sigma},
\end{equation}
where $X_{[\mu\nu]}:=X_{\mu\nu}-X_{\nu\mu}$; and with $\tilde g=\Omega_{1/2}^2g$,
$\Omega_{1/2}^2=1+\frac{R}{4M^2}-\frac1{M^2}\bigl(\frac{\square\sqrt\sigma}{\sqrt\sigma}-\frac\Lambda\sigma\bigr)$
and $\tilde u_\mu=P_\mu/M$,
\begin{equation}\label{eq:spinhalflorentz}
\boxed{M\,\tilde u^\nu\tilde\nabla_\nu\tilde u_\mu=\tilde u^\nu F_{\nu\mu}.}
\end{equation}
\end{theorem}
\noindent The convective flow moves in the conformally rescaled metric like a
charge in the spin vorticity. $F$ contains a curvature--spin term with the
structure of the Mathisson--Papapetrou coupling, derived here rather than
imported. For spin-0 matter $\Lambda=0$ and $F=0$, and the law reduces to
Proposition~\ref{prop:qcf}.

\noindent\textbf{Verification.} Flat-region identities verified to $10^{-14}$ on exact
superpositions of plane-wave solutions (QGD--112). The curvature coefficients $+\frac14$
and $-\frac14$, the curved off-shell identity and the curved spin vorticity were verified
on an FLRW metric with an arbitrary spinor field, residuals $\le6\times10^{-16}$, with
the spinor connection constructed from $\nabla\gamma=0$ (QGD--113). In generic states the
spin-vorticity force is $40$--$70\%$ of the gradient force.

\noindent\textbf{Slow-motion limit (QGD--113).} With $\hat s$ the spin direction,
$\Lambda/\sigma\to\frac14\partial\hat s\cdot\partial\hat s$ (from the exact Fubini--Study identity
for unit two-spinors) and $F_{ij}\to-\frac12\hat s\cdot(\partial_i\hat s\times\partial_j\hat s)$, a
Mermin--Ho-type relation, both with $O(v^2)$ convergence verified. The kinetic
energy per particle is
$p^2/2M-(\hbar^2/2M)\nabla^2\sqrt n/\sqrt n+(\hbar^2/8M)|\nabla\hat s|^2$. For constant spin
the law reduces exactly to Chapter~1's irrotational result
$\dot{\bm p}=-\nabla(U+Q)$ and to the fixed-spin quantum stress of
Section~\ref{sec:quantumsource}. When spin varies, a spin quantum potential and a
spin-vorticity force appear that those fixed-spin treatments omit.

\noindent\textbf{Closed form of the spin-$\frac12$ stress (QGD--123).} Using the macroscopic spinor
equation on $\psi$ and on $\bar\psi$, with $\gamma^\lambda\gamma^\mu=-\eta^{\lambda\mu}-i\Sigma^{\lambda\mu}$, the symmetric
(Belinfante) stress tensor in a flat region is exactly
\begin{equation}\label{eq:Tclosed}
T^{\mu\nu}=\frac{\sigma}{M}P^\mu P^\nu
+\frac1{2M}\bigl(\partial^\mu\!\sqrt\sigma\,\partial^\nu\!\sqrt\sigma-\sqrt\sigma\,\partial^\mu\partial^\nu\!\sqrt\sigma\bigr)
+\frac1M\Lambda^{\mu\nu}
-\frac1{4M}\partial_\lambda\operatorname{Im}\bigl(\bar\psi\Sigma^{\lambda\mu}\partial^\nu\psi+\bar\psi\Sigma^{\lambda\nu}\partial^\mu\psi\bigr),
\end{equation}
with the tensor spin remainder
$\Lambda^{\mu\nu}=\operatorname{Re}(\partial^\mu\bar\psi\partial^\nu\psi)-\operatorname{Re}(\bar C^\mu C^\nu)/\sigma$, whose trace is $\Lambda$. The terms
are: dust; a scalar quantum stress; a spin remainder; and a spin curl, which is separately
conserved because $\Sigma^{\lambda\mu}$ is antisymmetric. Verified on exact superpositions of
plane-wave solutions (relative residual $\le5\times10^{-13}$). Relative to dust, the scalar and
spin-remainder terms are second order in $\lambda_C/L$ (slopes $1.92$, $1.94$) and the spin curl is
first order (slope $1.02$). The spin curl is the exact origin of the first-order tilt between
probability flow and energy flow. With $\hbar$ restored, the quantum stress of the universal
level for spin-$\frac12$ matter is therefore given in closed form. The curved-metric version
(replacing $\partial$ by the spinor covariant derivative) adds curvature terms that have not been
derived here.

\noindent\textbf{Conditions.} $\sigma\neq0$ and $P$ timelike. As for spin-0 matter, the
law describes the convective momentum field; reading its integral curves as
trajectories is an interpretation.

\subsection{Where the natural invariant comes from}
\begin{proposition}[Closure identity (D)]\label{prop:closure}
On the static one-null-sector branch $g=\bar g+H\,\mathrm dv^2$ with $H\to0$ at
infinity, wherever $0<H<1$,
\begin{equation}\label{eq:closure}
\sigma_{\rm amp}=\sqrt H=\beta_{\rm ff}=\frac{2\zeta_{\rm ff}}{1+\zeta_{\rm ff}^2},\qquad
\zeta_{\rm ff}=\frac{\sqrt H}{1+\sqrt{1-H}},
\end{equation}
where $\beta_{\rm ff}$ is the speed, relative to static observers, of a test
particle released from rest at infinity, and $\zeta_{\rm ff}$ is its Chapter~1
on-shell spinor ratio. Covariantly, $\sigma_{\rm amp}^2=1-(u_{\rm ff}\cdot u_{\rm stat})^{-2}$.
The sector direction $k=\mathrm d(t+r)$ is a phase gradient solving the massless
limit of Chapter~1's Hamilton--Jacobi equation.
\end{proposition}
\noindent This extends Chapter~2's Schwarzschild diagnostic to the whole static
one-null-sector class. It fails for Kerr: the free-fall speed differs from the
Kerr--Schild amplitude, and the oblate $k$ is not closed. It also fails where
$FH\neq1$, for example in Eq.~\eqref{eq:ABquantum}. It establishes
self-consistency, not origin. Chapter~1 (Proposition~QGD--116) uses it to repair the
historical source profile: the historical $b_f=\sqrt f$ is the velocity, and the spinor
ratio is $\zeta_f=\sqrt f/(1+\sqrt{1-f})$, which matches exact radial solutions to
$3\times10^{-4}$. The natural invariant's algebraic \emph{form} is
derived from the macroscopic spinor equation, while its exterior \emph{values}
for a given source are selected by F4.

\subsection{Lorentz covariance and gauge dependence}
The main equations were tested directly (QGD--110):
\begin{center}\small
\begin{tabular}{@{}p{0.45\linewidth}p{0.47\linewidth}@{}}\toprule
Equation or object & Result\\\midrule
Master metric, Eq.~\eqref{eq:master} & Invariant under frame Lorentz transformations ($T\mapsto\Lambda T$, $\sigma_a\mapsto\sigma_c(\Lambda^{-1})^c{}_a$); tensorial under coordinate changes\\
Historical $\mathcal Q$ with $\bar\nabla$ & Tensorial: Cartesian and spherical computations agree exactly\\
Macroscopic spinor equation (Ch.~1 conventions) & Covariant: a spinor representation $S(\Lambda)$ exists and $S u(p)$ solves the equation at $\Lambda p$\\
Spinor ratio $\mathsf R(p)$ & Transforms exactly as $\mathsf R(p)\mapsto\mathsf R(\Lambda p)$\\
$\zeta=p/(E+Mc^2)$ & Poincar\'e-ball coordinate of the mass shell; boosts act by M\"obius addition (exact). $\beta=2\zeta/(1+\zeta^2)$ is the Beltrami--Klein coordinate\\
Spinor component moduli, $\psi^\dagger\psi$ & \textbf{Not} invariant; $\bar\psi\psi$ and $j^\mu j_\mu$ are\\
$\Omega_\psi$, $T^Q$ (scalar wave) & Covariant\\
Sector representation $\sigma$ & \textbf{Gauge dependent}: the same Schwarzschild geometry (Kretschmann $12r_s^2/r^6$ in both) has $h=g-\bar g$ of rank 1 in the advanced identification and rank 2 in the static one\\\bottomrule
\end{tabular}
\end{center}
\noindent Three consequences follow:
\begin{enumerate}
\item \textbf{Spin-$\frac12$ version of the quantum terms (QGD--111--113).} The scalar-wave
Eqs.~\eqref{eq:qcf} and~\eqref{eq:TQ} use component amplitudes, which are not
invariant for spinors. The covariant spin-$\frac12$ amplitude is $(-j\cdot j)^{1/4}$
(Fierz identity, verified). The per-wave level for spin-$\frac12$ matter is
Theorem~\ref{thm:spinhalf}. The universal level holds unchanged: in the energy frame,
the stress tensor of a coherent spin-$\frac12$ field is $\rho\,u_Lu_L+\Pi$ with
$\Pi=O((\lambda_C/L)^2)\rho$ (numerical scaling, slope 1.98), and the probability
current tilts relative to the energy flow at $O(\lambda_C/L)$ (slope 1.01) through the
spin current, which is not an extra stress.
\item \textbf{Geometric meaning of the spinor ratio.} The Chapter~1 ratio $\zeta$
is not a Lorentz vector. It is a hyperbolic-geometry coordinate (Poincar\'e
ball), and the closure identity~\eqref{eq:closure} is the ball-to-Klein map.
\item \textbf{The natural invariant is not a diffeomorphism invariant.}
Because the master metric contains a fixed reference $\bar g$, the sector
representation, including the number of sectors, their rank and their signs,
depends on how spacetime is identified with the reference. Only $g$, and
quantities built from it, are gauge invariant. Every QGD prediction under F4 is
computed from $g$, so no prediction is affected, but statements about $\sigma$
must state the identification used.
\end{enumerate}
% =================== END QGD-110 INSERTION ===================

\section{What the equation makes possible---and what it does not}
\label{sec:physics}
Equation~\eqref{eq:split} has an important explanatory role.
It expresses effective gravitational equilibrium as a balance of the
metric-mediated gravitational action against physical matter stress,
projected into the natural-invariant variables. The source is the
\emph{complete, conserved} stress of the physical system; a local
Dirac-component ratio or a single wave amplitude is not generally an
adequate universal source. The Dirac state may contribute through
basis-independent currents, energies and stress correlations.

But the action-derived equation also establishes a \emph{limit}:
with no direct natural-invariant interactions and no new dynamical
operator, the effective field equations contain no distinct
propagating $\sigma$ mode beyond the effective metric degrees of
freedom. A new observable could follow from (i) a well-defined direct
coupling between quantum matter and natural-invariant data, (ii) a
physically derived independent natural-invariant kinetic/constraint
sector, or (iii) a genuinely distinct quantum effective action.
None follows inevitably from the signed metric identity; each would
require a separately specified covariant action, conservation laws,
correct field content and measurable consequences.

% === QGD-101 insertion, corrected QGD-105 ===
\paragraph{Note on route (i) (QGD--100, corrected QGD--105).}
Research note QGD--100 examined one \emph{hypothetical} realization of route~(i):
removing the reference term, reading four $\eta$-signed sectors as a tetrad, and
assuming a microscopic theory with only global frame symmetry. Neither assumption
belongs to the formulation of Chapters~1--4. Under them the two-derivative
dynamics is New General Relativity \cite{HS}, whose parameter directions other than the
axial one give $\gamma=-2c/(a+b-4c)\ne1$ at linear order, and whose axial
direction is obstructed at nonlinear order according to \cite{BJD}.
This rules out that particular hypothetical route; it says nothing about
routes (i)--(iii) in general.
% === end QGD-101/105 insertion ===

A particularly sharp discriminator is the comparison of different
physical quantum source states with matched mean stress and classical
boundary data. A measured mean-field difference beyond standard
GR plus quantum-field corrections would indicate a new physical
coupling. Merely assigning different signed representatives to the
same metric cannot change a metric-only observable.

\section{Relationship to the archived historical Chapter 4}
\label{sec:historical}
The historical research program wrote schematically
\begin{equation}\label{eq:historical}
\square_g\sigma_\mu=Q_\mu+G_\mu+T_\mu
+\kappa\ell_Q^2\square_g^2\sigma_\mu+\cdots.
\end{equation}
This is \emph{not} equivalent by definition to
Eq.~\eqref{eq:split}. The historical $Q_\mu,G_\mu,T_\mu$
expressions have unresolved algebraic/provenance and normalization
issues documented in the QGD--052 audit. In particular, no term
$\square_g\sigma_\mu$ or $\square_g^2\sigma_\mu$
can be deduced simply from varying the derivative-free composite
Einstein--Hilbert action. Adding a standalone $\sigma$ kinetic term
would be a change of the effective theory and requires physical
justification, not a cosmetic rewriting of the balance equation.
The archived manuscript therefore remains an independent historical
source, while the present chapter begins with the precise equation
supported by the accepted action and metric map.

\section{Conclusion and the next legitimate question}
The action-based content of this updated Chapter~4 has one central field-equation result:
\begin{equation}\label{eq:conclusion}
\boxed{\cS_{(s)}^\rho
=2s_s\cE_g^{\rho\nu}\sigma_\nu^{(s)}=0,\qquad
\cS_{(s)}^\rho=\cG_{(s)}^\rho+\cT_{(s)}^\rho+\cQ_{(s)}^\rho.}
\end{equation}
For the metric-only effective action, $\cQ=0$ and, on the
full-rank signed branch, this balance reproduces the established
Einsteinian field equations. This is a rigorous mathematical bridge
from QGD's natural-invariant variables to effective gravity, not a
claim that the underlying physical origin of gravity has been solved.
The question that can make QGD more than a geometric reformulation
is therefore precise: \emph{Does the quantum matter-wave structure
lead to independent, covariantly defined natural-invariant dynamics
or source dependence, and does any such effect produce an observable
not already predicted by GR and quantum field theory?}
\par\medskip
The Newton-calibrated Chapter~3 potential and the controlled $\hbar^2$
quantum-source stress are compatible with this balance, but neither
supplies an independent natural-invariant propagation operator.
The metric-only quantum-equivalence theorem in
Section~\ref{sec:independent} makes this limitation exact under
its stated rank and action assumptions; the historical attractive
$9/2$ force would additionally require a nonzero exterior source
residual or a changed gravitational response. The
historical $9/2$ correction remains unverified. Future extensions must
derive their covariant operators, coefficient matching and physical
observables before assigning a genuinely new wave equation or degree
of freedom to $\sigma_\mu$.


% === QGD-101 insertion (moved from the abstract) ===
\section*{Revision history}\label{sec:revision-history}
\addcontentsline{toc}{section}{Revision history}
\textbf{QGD--096} established the equivalence/no-independent-correction
theorem for the regular metric-only theory and calculated the nonzero
vacuum-Poisson residual that the historical attractive $9/2$ force would
require. \textbf{QGD--097} added a separately sourced semiclassical
benchmark for a massless four-component Dirac vacuum field: a verified
far-field $-2G_N^2M\hbar/(15\pi c^3r^3)$ correction to the static
Killing-energy potential, taken from published quantum-matter backreaction,
not from the historical complex-phase Taylor expansion or a new QGD
operator. \textbf{QGD--098} independently derived and checked the
angular-channel coefficients of the Schwarzschild radial Dirac reduction and
recorded why its Liouville-normal effective potential cannot be promoted to
a physical force by differentiation. \textbf{QGD--101} integrated research
notes QGD--099 and QGD--100. \textbf{QGD--105} audited those insertions:
Remark~\ref{rem:Q-absorbed} was restricted to Chapter~2's constructed branch;
the route-(i) paragraph was restated as a test of one hypothetical realization;
and a closing paragraph that reframed QGD's identity was removed, because it
did not follow from the chapter's stated assumptions. No earlier equation or
label was removed. \textbf{QGD--110} added Section~\ref{sec:quantumstructure}: the
split identity, the quantum term in static variables, the test of the historical
derivative trial, the derived quantum conformal factor, the two-level master metric,
the closure identity, and Lorentz-covariance and gauge tests. \textbf{QGD--111} replaced
the open spin-$\frac12$ item with covariant results; \textbf{QGD--112} added the exact
spin-$\frac12$ flow law and restricted the name ``Dirac'' to the microscopic field.
\textbf{QGD--113} added the curved-metric form of that law (Theorem~\ref{thm:spinhalf}),
its slow-motion limits, and the cross-reference in Section~\ref{sec:quantumsource}.
\textbf{QGD--123} added the closed form of the spin-$\frac12$ stress, Eq.~\eqref{eq:Tclosed}, and
applied the naming convention to derived objects. \textbf{QGD--124} added the abstract sentence on
QGD's place as quantum gravity.
\textbf{QGD--120} added Section~\ref{sec:qgdaction}: the complete QGD action, its symmetries
and its quantum effective form. \textbf{QGD--121} added Section~\ref{sec:dof}: the degree-of-freedom
count (two propagating modes on spanning branches; degenerate and non-spanning cases).
\textbf{QGD--122} added Proposition~\ref{prop:backgrounds}: spanning representations of the
chapters' black-hole and cosmological backgrounds, and the projection form of one-sector reductions.
\textbf{QGD--115} added Section~\ref{sec:lovelock}: the Einstein--Hilbert form derived by
uniqueness under Principle~L, with the coupling from the Newtonian calibration.
% === end QGD-101 insertion ===

\section*{Source and research provenance}
\begin{thebibliography}{9}
\bibitem{ch1} QGD, \emph{Chapter 1: Foundations}, accepted Dirac/matter-wave construction (existing research manuscript).
\bibitem{ch2} QGD, \emph{Chapter 2: Field Equations}, complete signed master metric, variational right-inverse and effective action (QGD--075).
\bibitem{ch3} QGD, \emph{Chapter 3: Gravitational Normalization and Quantum-Force Verification} (QGD--087 surgical revision; QGD--101 notation update).
\bibitem{audit} QGD--052, \emph{Focused Historical QGT/QGD Field-Equation Provenance Audit}; and QGD--094, \emph{General Natural-Invariant Field Equations and Historical Source Split}, working research checks.
\bibitem{referee} QGD--091--092, \emph{Independent Schwarzschild Dirac radial reduction and referee verification}, conditional test-field coefficient audit (cross-checked in QGD--098).
\bibitem{quantum} QGD--090--094, \emph{Quantum Force, Dirac Radial Dynamics, and Quantum-Source Verification}, conditional source and effective-metric studies. Historical $\hbar^2$ operators remain unmatched.
\bibitem{Carlson2010} E. D. Carlson, P. R. Anderson, A. Fabbri, S. Fagnocchi, W. H. Hirsch, and S. A. Klyap, ``Semiclassical gravity in the far field limit of stars, black holes, and wormholes,'' \emph{Phys. Rev. D} \textbf{82}, 124070 (2010), DOI: 10.1103/PhysRevD.82.124070; see Eqs. (7.2), (8.3), (8.6), and (8.7).
\bibitem{donoghue} J. F. Donoghue, ``General Relativity as an Effective Field Theory: The Leading Quantum Corrections, \emph{Phys. Rev. D} \textbf{50}, 3874 (1994), DOI: 10.1103/PhysRevD.50.3874.
\bibitem{equivalence} S. Kamefuchi, L. O\textquoteright Raifeartaigh, and A. Salam, ``Change of variables and equivalence theorems in quantum field theories, \emph{Nuclear Physics} \textbf{28}, 529--549 (1961), DOI: 10.1016/0029-5582(61)91075-6.
\bibitem{qgd100} QGD--100, \emph{Testing the Global-Lorentz Branch}, research note (2026).
\bibitem{HS} K.~Hayashi and T.~Shirafuji, ``New general relativity,'' \emph{Phys. Rev. D} \textbf{19}, 3524 (1979).
\bibitem{Lovelock1971} D.~Lovelock, ``The Einstein tensor and its generalizations,'' \emph{J. Math. Phys.} \textbf{12}, 498 (1971).
\bibitem{Lovelock1972} D.~Lovelock, ``The four-dimensionality of space and the Einstein tensor,'' \emph{J. Math. Phys.} \textbf{13}, 874 (1972).
\bibitem{Jacobson} T.~Jacobson, ``Thermodynamics of spacetime: the Einstein equation of state,'' \emph{Phys. Rev. Lett.} \textbf{75}, 1260 (1995).
\bibitem{Sakharov} A.~D.~Sakharov, \emph{Sov. Phys. Dokl.} \textbf{12}, 1040 (1968).
\bibitem{Visser} M.~Visser, ``Sakharov's induced gravity: a modern perspective,'' \emph{Mod. Phys. Lett. A} \textbf{17}, 977 (2002).
\bibitem{BJD} J.~Beltr\'an Jim\'enez and K.~F.~Dialektopoulos, ``Non-linear obstructions for consistent new general relativity,'' \emph{JCAP} \textbf{2020}(01), 018, arXiv:1907.10038.
\end{thebibliography}

% ========================================================================
% QGD-138 / Chapter 5 v3.3: Physical quantum backreaction vs historic Q
% ========================================================================
\section{The quantum correction slot: what QGD actually derives}\label{sec:qgd138}
\noindent\textbf{Scope and proof status.} The complete founding master metric is
\begin{equation}\label{eq:qgd138-master}
 \boxed{g_{\mu\nu}=T^a{}_{\mu}T^b{}_{\nu}
 \biggl[\eta_{ab}-\sum_s\varepsilon_s\sigma_a^{(s)}\sigma_b^{(s)}
 -\lambda\mathcal Q_{ab}\biggr].}
\end{equation}
Here the historical $\mathcal Q_{ab}$ is a proposed independent
metric-derivative operator. We distinguish it from (i) a computed,
state-dependent quantum metric deformation $q_{ab}^{\rm back}$ and
(ii) the phase/amplitude quantum potential of an individual matter wave.
We prove a linear response identity and a metric-level non-identifiability
statement; neither constitutes a quantum microscopic derivation of a
unique new $\mathcal Q[\sigma]$ operator.

\subsection{The response that follows from the renormalized action}
Let $g^{(0)}$ solve the selected classical effective Einstein--matter
problem and assume a controlled perturbative state with
$\langle\widehat T_{\mu\nu}\rangle_{\rm ren}=O(\hbar)$.
For $g=g^{(0)}+q^{\rm back}+O(\hbar^2)$, expansion of the
quantum-corrected metric equations at first order yields
\begin{equation}\label{eq:qgd138-eom}
 \boxed{\mathcal L_{g^{(0)}}[q^{\rm back}]_{\mu\nu}
 =\frac{8\pi G_N}{c^4}\langle\widehat T_{\mu\nu}\rangle_{\rm ren}^{(1)}.}
\end{equation}
The linear operator $\mathcal L_{g^{(0)}}$ includes the linearized
classical geometric and matter Euler derivatives; if quantum
polarization is resummed, the corresponding causal stress response
kernel must also be included. In a chosen gauge and with admissible
boundary/initial data, the formal first-order solution is
\begin{equation}\label{eq:qgd138-green}
 \boxed{q^{\rm back}=\frac{8\pi G_N}{c^4}
 \mathcal L_{g^{(0)},\,\rm bc}^{-1}
 \langle\widehat T\rangle_{\rm ren}^{(1)}.}
\end{equation}
For real-time evolution the inverse is a retarded, gauge-constrained
response, not an arbitrary Feynman Green function. This is a
\emph{derived quantum metric response conditional on the Chapter 4 action},
not a new off-shell operator in the defining metric.

\begin{proposition}[On-shell historical-slot assignment, convention-dependent]\label{prop:qgd138-slot}
Fix a reference coframe, a gauge/coordinate identification, a classical
natural-invariant representative, a quantum state and boundary data.
Choose $\lambda\ne0$. If the sector configuration is held fixed while
representing the first-order quantum correction, one may define
\begin{equation}\label{eq:qgd138-slot}
 \boxed{\mathcal Q_{ab}^{\rm rep}:=-\frac1\lambda
 (T^{-1})^\mu{}_a(T^{-1})^\nu{}_b\,
 q_{\mu\nu}^{\rm back}.}
\end{equation}
Then $-\lambda T^a{}_{\mu}T^b{}_{\nu}
\mathcal Q_{ab}^{\rm rep}=q_{\mu\nu}^{\rm back}$ by construction.
This is an \emph{on-shell representation} of a computed metric
correction; it is not a unique or autonomous functional
$\mathcal Q_{ab}[\sigma,\bar\nabla\sigma]$ and does not
supply the adjoint $D_s\mathcal Q$ for a new field equation.
\end{proposition}

\subsection{A calculated far-field example: one massless Dirac field}
Use the Boulware-state far-field solution and the static, fixed-areal-radius
coordinates of Chapter 4. Define the length scales
$\mu:=G_NM/c^2$ and $\ell_P^2:=G_N\hbar/c^3$, and
$f_0:=1-2\mu/r$. The one-loop massless-Dirac solution gives, at
leading quantum order for $r\gg2\mu$,
\begin{equation}\label{eq:qgd138-correctedf}
 f(r)=f_0-\frac{4\mu\ell_P^2}{15\pi r^3}+\cdots,
 \qquad
 f_r(r):=1-\frac{2G_N\mathfrak m(r)}{c^2r}
 =f_0-\frac{2\mu\ell_P^2}{5\pi r^3}+\cdots .
\end{equation}
The common-coordinate metric difference relative to classical
Schwarzschild is therefore
\begin{equation}\label{eq:qgd138-qcomponents}
 \boxed{q^{\rm back}_{00}=\frac{4\mu\ell_P^2}{15\pi r^3},
 \qquad q^{\rm back}_{rr}=\frac{2\mu\ell_P^2}{5\pi r^3}
 +O\!\left(\frac{\mu^2\ell_P^2}{r^4}\right),\qquad
 q^{\rm back}_{\theta\theta}=q^{\rm back}_{\phi\phi}=0.}
\end{equation}
Here $x^0=ct$ and the covariant metric components are used; the
angular components vanish in areal gauge. Other choices of gauge change
components, not physical curvature observables.
Selecting, only as a bookkeeping convention, $\lambda=\ell_P^2$,
and the diagonal flat spherical reference coframe,
Eq.~\eqref{eq:qgd138-slot} yields the explicit representative
\begin{equation}\label{eq:qgd138-Qexample}
 \boxed{\mathcal Q_{00}^{\rm rep}=-\frac{4\mu}{15\pi r^3},
 \qquad \mathcal Q_{11}^{\rm rep}=-\frac{2\mu}{5\pi r^3}
 +O\!\left(\frac{\mu^2}{r^4}\right),\qquad
 \mathcal Q_{22}^{\rm rep}=\mathcal Q_{33}^{\rm rep}=0.}
\end{equation}
The representative has the required dimension $L^{-2}$ if the
natural invariants are dimensionless and $\lambda$ has dimension $L^2$.
It vanishes from the metric as $\hbar\to0$, provided the state and
classical background have a regular classical limit. No coefficient
of an independent new QGD force has been obtained.

\begin{proposition}[No unique $\mathcal Q$ from metric-only dynamics]\label{prop:qgd138-nounique}
At any point with a nondegenerate reference coframe, the complete real
signed-sector construction can represent an arbitrary symmetric metric
deformation. Therefore a specified total metric and a metric-only
effective action do not select a unique additional tensor
$\mathcal Q_{ab}$. Even with the quantum metric $q^{\rm back}$ computed,
one may use an on-shell representative with the historical slot
$-\lambda\mathcal Q=q^{\rm back}$ and fixed classical sectors,
\emph{or} readjust sufficiently many signed sectors and take
$\mathcal Q=0$, obtaining the same metric.
\end{proposition}
\begin{proof}
Write the total internal deformation as
$h_{ab}=(T^{-1}gT^{-\mathsf T})_{ab}-\eta_{ab}$.
For any symmetric $h$, choose $b\ne0$ and the signed matrices
$P=bI+h/(4b)$ and $N=bI-h/(4b)$. They obey
$P^{\mathsf T}P-N^{\mathsf T}N=h$ exactly. Their rows supply a
real signed-sector representation with $\mathcal Q=0$.
Alternatively, retain a chosen classical set of sectors and assign
$-\lambda\mathcal Q$ to the remaining deformation. Because the
metric-only action sees only the resulting $g$, it does not
discriminate these metric representations. The argument is
algebraic and on-shell: it does not establish equivalence of all
\emph{off-shell} derivative-dependent changes of fields or their
quantum functional measures.
\end{proof}

\subsection{Why a historical local derivative candidate is not enough}
The trial $\mathcal Q_{ab}^{\rm trial}=
(\bar\nabla_a\sigma^c)(\bar\nabla_b\sigma_c)$
has the correct schematic dimensional order: on the Schwarzschild
background $\sigma\sim\sqrt{\mu/r}$ and
$\ell_P^2\mathcal Q^{\rm trial}\sim\mu\ell_P^2/r^3$.
But Chapter 4's independent tensor test found that no one coefficient
reproduces \emph{both} the quantum temporal and radial metric corrections.
Matching a radial power law does not match a covariant tensor.
The known renormalized matter effective action also contains
state-dependent and generally nonlocal responses, which a single
local first-gradient ansatz need not reproduce.

\paragraph{Consequence for Option C.}
For a genuinely specified differentiable off-shell tensor
$\mathcal Q[\sigma;T]$, the preferred QGD equation is
\begin{equation}\label{eq:qgd138-optionC}
 \boxed{2s_s\Pi^{ab}\sigma_b^{(s)}=
 \lambda\bigl[(D_s\mathcal Q)^*_{\sqrt{-g}}\Pi\bigr]^a,\qquad
 s_s=-\varepsilon_s.}
\end{equation}
But substituting the on-shell representative
\eqref{eq:qgd138-slot} into this equation as if it were an
independent operator would generally double count or misrepresent the
same quantum backreaction already present in
$\hbar\Gamma_{\rm matter}^{(1),\rm ren}$.
The presently controlled dynamical formulation sets the independent
historical slot to zero and retains its quantum source in $\Pi^{ab}$.
A nonzero new $\mathcal Q$ requires an independently defined
sector-sensitive observable, a microscopic matching calculation, or a
nonredundant dynamical principle.
% ========================================================================


% ===== QGD-139: Energy-charge and field-equation interpretation; 31 October 2026 =====
\section{Energy charges and the meaning of the projected field equations (QGD--139)}
\label{sec:qgd139-energycharges}
\textbf{Premises and scope.} This section uses the complete founding metric
$g_{\mu\nu}=T^a{}_{\mu}T^b{}_{\nu}[\eta_{ab}-\sum_s\varepsilon_s\sigma_a^{(s)}\sigma_b^{(s)}-\lambda\mathcal Q_{ab}]$.
The explicit calculations are on its established fixed-coframe,
$\lambda\mathcal Q=0$ metric-only branch. The historical $\mathcal Q$
remains uncalculated; a background translation current is not thereby
proved to be a unique gauge-invariant local gravitational stress tensor.

\subsection{Two levels of energy conservation}
For Chapter~1's flat-space matter action, translation invariance implies
$\partial_\mu\Theta^\mu{}_{\nu}=0$ on the microscopic matter equations.
For Chapter~4's emergent effective metric, the covariant Bianchi identity
and metric matter equations give $\nabla_\mu T_{\rm eff}^{\mu\nu}=0$,
including the renormalized quantum source in its regime of validity.
They are not the same local current. Where $\xi^\nu$ is a Killing field,
$J^\mu:=T_{\rm eff}^{\mu\nu}\xi_\nu$ obeys
\begin{equation}\label{eq:qgd139-killing}
\nabla_\mu J^\mu=0,\qquad
\partial_\mu(\sqrt{-g}\,J^\mu)=0.
\end{equation}
This is a matter Killing-energy balance; in a gravitating system the
Hamiltonian energy additionally includes a gravitational boundary term.
For nonstationary systems an asymptotic, e.g. Bondi, flux replaces a
global stationary Killing energy.

\subsection{The Option C equation and the conservation constraints}
For the metric-only functional, define $\Pi^{ab}:=T^a{}_{\mu}T^b{}_{\nu}\mathcal E_g^{\mu\nu}$,
where $\mathcal E_g^{\mu\nu}=(\sqrt{-g})^{-1}\delta\Gamma/\delta g_{\mu\nu}$.
The complete metric-map projection is
\begin{equation}\label{eq:qgd139-optionC}
\boxed{2s_s\Pi^{ab}\sigma_b^{(s)}=
\lambda[(D_s\mathcal Q)^*_{\sqrt{-g}}\Pi]^a,\qquad s_s=-\varepsilon_s.}
\end{equation}
On a regular branch with a surjective complete metric Jacobian,
Eq.~\eqref{eq:qgd139-optionC} is equivalent to $\Pi=0$, that is the
quantum-corrected Einstein--matter equations. Thus both displayed sides
vanish on shell: they are \emph{not} separately observable opposing
energy-transfer channels. The energy conservation law is instead obtained
from the diffeomorphism/translation Noether identities and boundary
Hamiltonian structure. A new physical sector energy requires a separately
justified sector-sensitive action, which the established theory lacks.

\subsection{Schwarzschild: matching the natural-invariant flux to ADM mass}
Let $H=A^2=B^2=2G_NM/(c^2r)=C/r$ for the standard flat-reference
Schwarzschild Kerr--Schild exterior. On a $t=\mathrm{const}$ Kerr--Schild
slice in asymptotically Cartesian coordinates, $\gamma_{ij}=\delta_{ij}+Hn_in_j$,
$n_i=x_i/r$. Since
\begin{equation}\label{eq:qgd139-adm-integrand}
\partial_j(\gamma_{ij}-\delta_{ij})
-\partial_i(\gamma_{jj}-3)=\frac{2H}{r}\,n_i
\qquad\text{for }H=C/r,
\end{equation}
the ADM mass is
\begin{equation}\label{eq:qgd139-adm}
\boxed{M_{\rm ADM}=\frac{c^2}{16\pi G_N}
\lim_{r\to\infty}\oint
\bigl[\partial_j\gamma_{ij}-\partial_i\gamma_{jj}\bigr]n^i\,dS
=\frac{c^2 C}{2G_N}=M.}
\end{equation}
Independently, Chapter~5's accepted static energy gives
$\nabla^2H=-8\pi G_N\rho_m/c^2$ and hence its source flux
\begin{equation}\label{eq:qgd139-flux}
\boxed{M_{\rm flux}:=-\frac{c^2}{8\pi G_N}
\oint \partial_rH\,dS=\frac{c^2C}{2G_N}=M_{\rm ADM}.}
\end{equation}
This is an exact \emph{exterior-charge correspondence} for the named
Kerr--Schild branch and asymptotic conditions; identifying the Newtonian
source integral with a fully relativistic interior ADM mass requires the
appropriate internal, kinetic and binding contributions, not simply
$\int\rho_m\,d^3x$ beyond the Newtonian order.

\subsection{Exact spherical mass as a function of the natural invariants}
\label{sec:qgd139-ms}
For Chapter~2's one-sector spherically symmetric metric in an areal
radius $r$, put $F=1-A^2$ and $D=1+B^2-A^2$. The inverse radial
metric coefficient is $g^{rr}=F/D$, independent of the
static off-diagonal $dt\,dr$ representation. The standard
Misner--Sharp mass therefore has the exact QGD representation
\begin{equation}\label{eq:qgd139-misnersharp}
\boxed{m_{\rm MS}(r)=\frac{c^2r}{2G_N}
\bigl(1-g^{ab}\partial_a r\partial_b r\bigr)
=\frac{c^2r}{2G_N}\frac{B(r)^2}{1+B(r)^2-A(r)^2}.}
\end{equation}
This is an identity in the specified spherically symmetric metric
representative, not a new GR mass or a unique local stress tensor.
On $A^2=B^2=H=2G_NM/(c^2r)$ it gives $m_{\rm MS}=M$
identically in the vacuum exterior. In a static Einstein--matter
solution its radial derivative satisfies the geometric-source
balance $m_{\rm MS}'(r)=4\pi r^2\epsilon(r)/c^2$ for total local
rest-frame energy density $\epsilon$ in the conventional static
spherical frame, including renormalized quantum contributions
when these are consistently included.

In Chapter~4's specific massless-Dirac asymptotic calculation,
$A^2=H+\delta A^2$, $B^2=H+\delta B^2$,
where $\delta A^2=4G_N^2M\hbar/(15\pi c^5r^3)$ and
$\delta B^2=2G_N^2M\hbar/(5\pi c^5r^3)$. First-order expansion gives
\begin{equation}\label{eq:qgd139-msquantum}
\boxed{\delta m_{\rm MS}=
\frac{c^2r}{2G_N}\left[(1-H)\delta B^2+H\delta A^2\right]
=\frac{G_NM\hbar}{5\pi c^3r^2}
+O\!\left(\frac{G_N^2M^2\hbar}{c^5r^3}\right).}
\end{equation}
It reproduces the independently calculated radial mass function of
Chapter~4 (geometrized $m(r)=M[1+1/(5\pi r^2)]+\cdots$).
The exterior quantum stress can make this finite-radius mass larger
than its fixed asymptotic ADM value; $m_{\rm MS}'$ is then negative
in that state. This is a check of the effective equations, not a
new source-dependent quantum force from sector reparameterization.

\subsection{An exact static-spherical QGD to Brown--York energy dictionary}
\label{sec:qgd139-by-general}
For Chapter~2's general one-sector static spherical metric, write
$\sigma=A(r)\,d(ct)+B(r)\,dr$ and define
$F:=1-A^2$, $D:=1+B^2-A^2$. After removal of the static cross term,
the radial component on the orthogonal static slice is
$\gamma_{rr}=D/F$. The flat-referenced Brown--York energy of a
round sphere with areal radius $R$ therefore becomes
\begin{equation}\label{eq:qgd139-byAB}
\boxed{E_{\rm BY}(R)=\frac{c^4R}{G_N}\left[
1-\sqrt{\frac{1-A(R)^2}{1+B(R)^2-A(R)^2}}\right],
\quad F,D>0.}
\end{equation}
\emph{Proof:} On this spatial slice, the outward mean curvature is
$k=(2/R)\sqrt{F/D}$; the flat reference has $k_0=2/R$.
The Brown--York Hamiltonian surface integral
$E=(c^4/8\pi G_N)\oint(k_0-k)\,dS$
reduces to Eq.~\eqref{eq:qgd139-byAB}. For the exterior classical
branch $A^2=B^2=H$ this reduces to Eq.~\eqref{eq:qgd139-by}.
It is a quasilocal mass in the specified static slice and reference;
for other slicings and observers the Brown--York energy differs.

For the massless-Dirac far-field backreaction, let
$A^2=H+\delta A^2$ and $B^2=H+\delta B^2$ with
$H=2G_NM/(c^2R)$ and
$\delta A^2=4G_N^2M\hbar/(15\pi c^5R^3)$,
$\delta B^2=2G_N^2M\hbar/(5\pi c^5R^3)$.
Differentiating Eq.~\eqref{eq:qgd139-byAB} to first order in
$\delta A^2,\delta B^2$ gives
\begin{equation}\label{eq:qgd139-byQ}
\boxed{\delta E_{\rm BY}(R)=\frac{c^4R}{2G_N}
\sqrt{1-H}\left[\delta B^2+\frac{H}{1-H}\delta A^2\right]
=\frac{G_NM\hbar}{5\pi cR^2}
+O\!\left(\frac{G_N^2M^2\hbar}{c^3R^3}\right).}
\end{equation}
The equality of the leading term uses the specified Boulware state,
large $R$, and fixed asymptotic ADM mass, and is a reconstruction
of a known semiclassical metric rather than a distinct quantum
prediction of the signed representation. The last $O$ term denotes
higher classical-field powers multiplying this one-loop contribution.

\paragraph{Quantum backreaction and the asymptotic charge.}
Chapter~4's specified massless-Dirac quantum corrections change the
exterior radial and temporal metric amplitudes at order
$\delta A^2,\delta B^2=O(\hbar/r^3)$ with fixed asymptotic ADM
mass $M$. Their first spatial derivatives decay as $O(\hbar/r^4)$,
so they contribute $O(\hbar/r^2)$ to an ADM surface integral and
therefore no shift to its limit at infinity. Quantum matter still
changes finite-radius gravitational observables and energies.
This conclusion is specific to the named fixed-mass asymptotic state;
a different quantum state may also change the global mass integration
constant through its physical source and boundary conditions.

\subsection{The first quasilocal correction and its meaning}
Chapter~6's positive exterior static gradient contribution is
$E_{\rm grad}(r>R)=G_NM^2/(2R)$, while the same Schwarzschild geometry
on its conventional static spatial slice has the flat-referenced
Brown--York quasilocal energy
\begin{equation}\label{eq:qgd139-by}
E_{\rm BY}(R)=\frac{c^4 R}{G_N}
\left[1-\sqrt{1-\frac{2G_NM}{c^2R}}\right]
=Mc^2+\frac{G_NM^2}{2R}
+O\!\left(\frac{G_N^2M^3}{c^2R^2}\right).
\end{equation}
Therefore $E_{\rm BY}(R)-Mc^2=E_{\rm grad}(r>R)+O(R^{-2})$.
This is a \emph{leading weak-field coefficient match}, not equality
of local density or of quasilocal prescriptions: the Brown--York quantity
uses a specified reference and spatial slicing, while $E_{\rm grad}$
is part of an auxiliary static functional. Binding energy is negative
only after the matter--field interaction is included.

\subsection{Consequence for the interpretation of QGD}
In any purely metric-dependent QGD sector on which the metric map is
regular and boundary-compatible, the covariant gravitational
Hamiltonian/Noether surface charges are pulled back from those of the
same metric theory. ADM and Bondi charges therefore agree with the
metric comparator when the same physical geometry and boundary data
are used. The signed representation alone cannot remove GR's general
localization ambiguity; this is an exact limitation of the presently
accepted metric-only framework, not an obstacle to using its valid
energy balance and flux formulae. The reference coframe may help
organize calculations, but a unique new local energy requires a
physically selected frame/observer and a justified invariant prescription
(or new sector dynamics).

\noindent\textbf{Relevant references:}
Brown--York, \emph{Phys. Rev. D} 47 (1993) 1407,
\url{https://doi.org/10.1103/PhysRevD.47.1407};
Iyer--Wald, \emph{Phys. Rev. D} 50 (1994) 846,
\url{https://doi.org/10.1103/PhysRevD.50.846};
Szabados, \emph{Living Reviews in Relativity} 7 (2004),
\url{https://doi.org/10.12942/lrr-2004-4}.
% ===== END QGD-139 =====


% ===== QGD-140: Controlled quantum energy hierarchy; no new dynamics =====
\section{Quantum energy hierarchy and the limits of the natural-invariant energy dictionary (QGD--140)}
\label{sec:qgd140-quantumenergy}
\textbf{Scope.} Retain the complete founding master metric,
$g_{\mu\nu}=T^a{}_\mu T^b{}_\nu[\eta_{ab}-\sum_s\varepsilon_s\sigma_a^{(s)}\sigma_b^{(s)}-\lambda\mathcal Q_{ab}]$,
and the preferred Option C variational equation. Calculations here use the
\emph{previously established} metric-only, fixed-reference, $\lambda\mathcal Q=0$
semiclassical branch. No independent quantum energy of the historical
$\mathcal Q_{ab}$ is claimed or added to the metric effective action.

\subsection{Two independent perturbative parameters}
For an asymptotically flat source of fixed ADM mass $M$, set
\begin{equation}\label{eq:qgd140-scales}
 u(r):=\frac{G_NM}{c^2r},\qquad
 w(r):=\frac{G_N\hbar}{c^3r^2}=\frac{\ell_P^2}{r^2}.
\end{equation}
Terms of order $u^n$ can be purely classical; the source-specific Dirac
vacuum polarization below is of order $w$ relative to the leading
Newtonian force. Mixed $u w$ corrections are not automatically fixed by
knowing the leading far-field $w$ term.

\subsection{Massless Dirac vacuum energy in a fixed exterior state}
Adopt the specific renormalized asymptotically empty (Boulware)
state of one free massless four-component Dirac field and the
far-field semiclassical metric in Section~\ref{sec:qgd139-energycharges}.
The existing Chapter~4 calculation yields the leading areal-radius
Misner--Sharp mass, with $G_N,c,\hbar$ restored,
\begin{equation}\label{eq:qgd140-mass}
 m_{\rm MS}(r)=M+\delta m_1(r)+\cdots,
 \qquad \delta m_1(r)=\frac{G_N M\hbar}{5\pi c^3r^2}.
\end{equation}
Using the spherical Einstein constraint
$\dd m/\dd r=4\pi r^2\rho_E/c^2$ (where $\rho_E$ is local
static-frame energy density), its \emph{renormalized quantum matter}
contribution is
\begin{equation}\label{eq:qgd140-rho}
 \boxed{\rho_E^{(Q)}(r)=\frac{c^2}{4\pi r^2}\frac{\dd\delta m_1}{\dd r}
 =-\frac{G_NM\hbar}{10\pi^2c r^5}+\cdots.}
\end{equation}
Integrating the leading exterior tail gives the quantum matter
energy outside an areal sphere of radius $R$:
\begin{equation}\label{eq:qgd140-Eout}
 \boxed{E_{Q}^{\rm out}(R)
 =\int_R^\infty4\pi r^2\rho_E^{(Q)}(r)\,\dd r
 =-\frac{G_NM\hbar}{5\pi cR^2}+\cdots
 =-c^2\delta m_1(R)+\cdots.}
\end{equation}
The leading quantum matter energy is negative in this \emph{specified}
state. This is not a negative energy density assigned to a fundamental
$\sigma$ condensate or an assertion about arbitrary quantum states. At
fixed ADM mass $M$, the finite-radius enclosed mass is correspondingly
larger: $c^2[m_{\rm MS}(R)-M]=-E_Q^{\rm out}(R)$ at this order.

\subsection{Quasilocal energy, redshift and the first formal second-order term}
On the static-orthogonal, flat-subtracted Brown--York convention of
Chapter~6, let
\begin{equation}\label{eq:qgd140-BYfunctional}
 E_{\rm BY}(R;m)=\frac{c^4R}{G_N}
 \left[1-\sqrt{1-\frac{2G_N m(R)}{c^2R}}\right],
 \qquad f_0(R)=1-\frac{2G_NM}{c^2R}>0.
\end{equation}
Treat $m=M+\delta m_1+\delta m_2+\ldots$ with
$\delta m_n=O(\hbar^n)$ as an \emph{order-counting definition}.
Taylor expansion about $M$ gives the exact algebraic identities
\begin{align}
 \delta E_{\rm BY}^{(1)}(R)
 &=\frac{c^2\delta m_1(R)}{\sqrt{f_0(R)}},
 \label{eq:qgd140-by1}\\
 \delta E_{\rm BY}^{(2)}(R)
 &=\frac{c^2\delta m_2(R)}{\sqrt{f_0(R)}}
 +\frac{G_N[\delta m_1(R)]^2}{2R f_0(R)^{3/2}}.
 \label{eq:qgd140-by2}
\end{align}
The leading known term is therefore
\begin{equation}\label{eq:qgd140-BYknown}
 \delta E_{\rm BY}^{(1)}(R)
 =\frac{G_NM\hbar}{5\pi cR^2}
 +\text{uncomputed subleading one-loop terms}.
\end{equation}
The $f_0^{-1/2}$ factor in Eq.~\eqref{eq:qgd140-by1} is a legitimate
\emph{kinematical} response to a specified $\delta m_1(R)$; expanding
that factor does not determine the next far-field physical coefficient,
because omitted $O(\hbar r^{-3})$ contributions to $\delta m_1$ contribute
at the same order. Likewise the explicitly positive term quadratic in
$\delta m_1$ in Eq.~\eqref{eq:qgd140-by2} is only part of the
physical order-$\hbar^2$ correction: $\delta m_2$ remains unknown.
Obtaining $\delta m_2$ requires two-loop matter input and/or the
self-consistent first-loop response to the perturbed geometry,
including metric nonlinearities, counterterms and state conditions.

\subsection{A genuine second-order field-equation recursion}
For clarity use the expansion
$g=g_0+\hbar h_1+\hbar^2h_2+\cdots$,
$T_q[g]=\hbar\tau_1[g]+\hbar^2\tau_2[g]+\cdots$.
Let $\mathcal F[g]$ be the Einstein--classical-matter residual,
$\mathcal F[g_0]=0$, and $\kappa=8\pi G_N/c^4$. With
$\mathcal L:=D\mathcal F[g_0]$, expansion gives
\begin{align}
 \mathcal L[h_1]&=\kappa\tau_1[g_0],
 \label{eq:qgd140-h1}\\
 \mathcal L[h_2]&=\kappa\bigl(\tau_2[g_0]
 +D\tau_1[g_0]\cdot h_1\bigr)
 -\tfrac12 D^2\mathcal F[g_0](h_1,h_1).
 \label{eq:qgd140-h2}
\end{align}
These are perturbative identities under the stated expansion, not a
calculation of $\tau_2$. They distinguish the induced
$\hbar^2$ backreaction from a unique new fundamental QGD operator.
For genuinely quantum gravitational loops one must additionally
specify and quantize the relevant physical gravitational variables;
the current accepted action includes matter loops and does not
supply such a calculation.

\subsection{Meaning for QGD and gravitational energy}
The same $\sigma\mapsto g$ representation maps a corrected metric to
finite-radius energy observables, but the \emph{metric-only} QGD
field equations have the same ADM, Misner--Sharp and Brown--York
observables as their semiclassical GR counterparts in matched gauges,
states and boundary conditions. The projected Option C terms are not
individually conserved energy channels. The static Chapter~5 gradient
energy $G_NM^2/(2R)$ is positive, the on-shell two-body binding energy
negative, and the present exterior Dirac vacuum energy negative: they
refer to \emph{different functionals and sources} and therefore do not
contradict one another. A distinct QGD energy prediction demands a
new physically specified sector-sensitive observable or dynamics,
not an arbitrary further power of $1/r$ or $\hbar$.
\paragraph{A QGD-specific charge-matching question.}
The proposed flat-space microscopic starting point has a Noether
Hamiltonian $H_{\rm micro}$, while an asymptotically flat isolated
solution of the effective metric theory has an ADM Hamiltonian
$H_{\rm ADM}=M_{\rm ADM}c^2$. A microscopic collective-phase model
could be required to establish the \emph{conditional matching}
$\langle H_{\rm micro}\rangle_{\rm physical,ren}
\overset{?}{=}H_{\rm ADM}$, with the same reference energy, coarse-graining,
source and boundary prescription. This is a proposed test of QGD's
physical interpretation, not an identity proved by signed factorization
or by the Newtonian flux result.

\noindent\textbf{External benchmarks:} E. D. Carlson et al.,
\emph{Phys. Rev. D} \textbf{82}, 124070 (2010),
\url{https://doi.org/10.1103/PhysRevD.82.124070};
J. D. Brown and J. W. York, \emph{Phys. Rev. D} \textbf{47}, 1407 (1993),
\url{https://doi.org/10.1103/PhysRevD.47.1407};
J. F. Donoghue, \emph{Phys. Rev. D} \textbf{50}, 3874 (1994),
\url{https://doi.org/10.1103/PhysRevD.50.3874}.
% ===== END QGD-140 =====


% ===== QGD-141: historical source-balance re-derivation =====
\section{The historical three-channel equation revisited (QGD--141)}\label{sec:qgd141}
\textbf{Purpose and provenance.} The archival field-equations chapter proposed
$S_\mu=Q_\mu+G_\mu+T_\mu$ with $S_\mu=\Box_g\sigma_\mu$, and also
included trial higher-derivative terms. QGD--052 demonstrated that the
historical $Q_\mu$ identity, $T_\mu$ coefficient, and general derivation
of $G_\mu$ do not follow from the stated action. We therefore do not
revive $\Box_g\sigma_\mu$ as an established independent operator.
The architecture survives in an \emph{action-defined} successor, now
including the renormalized quantum matter functional.

\subsection{Exact three-channel construction using the full master metric}
Use the \emph{complete} founding map (Eq.~\ref{eq:qgd137-master}),
$g_{\mu\nu}=T^a{}_{\mu}T^b{}_{\nu}[\eta_{ab}-\sum_s
\varepsilon_s\sigma_a^{(s)}\sigma_b^{(s)}-\lambda\mathcal Q_{ab}]$,
with a fixed reference coframe and a specified differentiable
$\mathcal Q_{ab}[\sigma;T]$ when $\lambda\ne0$.
Separate the adopted renormalized metric-only action, after fixing a
renormalization scheme, as
\begin{equation}\label{eq:qgd141-splitaction}
\Gamma_{\mathrm{metric}}^{\mathrm{ren}}=
\underbrace{(S_{\mathrm{EH}}+S_{\mathrm{GHY}})}_{\Gamma_{\mathsf G}}
+\underbrace{S_{\mathrm{F1}}^{\mathrm{cl}}}_{\Gamma_{\mathsf T}}
+\underbrace{\bigl(\hbar\Gamma_{\mathrm{matter}}^{(1),\mathrm{ren}}
+O(\hbar^2)\bigr)}_{\Gamma_{\mathsf Q}},
\qquad \Delta\Gamma_{\mathrm{sector}}=0.
\end{equation}
For $X\in\{\mathsf G,\mathsf T,\mathsf Q\}$, define
$\Pi_X^{ab}:=T^a{}_{\mu}T^b{}_{\nu}
(\sqrt{-g})^{-1}\delta\Gamma_X/\delta g_{\mu\nu}$.
Then each channel has an exact full-map Euler derivative
\begin{equation}\label{eq:qgd141-channels}
\boxed{\mathscr X_{(s)}^{\,a}:=\frac{1}{\sqrt{-g}}
\frac{\delta\Gamma_X[g(\sigma)]}{\delta\sigma_a^{(s)}}
=-2\varepsilon_s\Pi_X^{ab}\sigma_b^{(s)}
-\lambda\bigl[(D_s\mathcal Q)^*_{\sqrt{-g}}\Pi_X\bigr]^a.}
\end{equation}
Linearity of functional differentiation proves the exact
\emph{stationarity balance}
\begin{equation}\label{eq:qgd141-balance}
\boxed{\mathscr S_{(s)}^{\,a}
:=\frac{1}{\sqrt{-g}}\frac{\delta\Gamma_{\mathrm{metric}}^{\mathrm{ren}}}
{\delta\sigma_a^{(s)}}
=\mathscr G_{(s)}^{\,a}+\mathscr T_{(s)}^{\,a}
+\mathscr Q_{(s)}^{\,a}=0.}
\end{equation}
Unlike historical $S_\mu=\Box_g\sigma_\mu$, this is an
\emph{Euler-residual equation} and not, without additional dynamics,
a free covector wave equation. The quantum channel
$\mathscr Q_{(s)}$ means the \emph{one-loop/higher-loop matter action
contribution}; it must \emph{not} be conflated with the separate
historical metric operator $\mathcal Q_{ab}$. This is an explicitly
\emph{updated naming convention}: the earlier QGD--098 decomposition
reserved $\mathscr Q$ for a derivative-map contribution and set
it to zero in the algebraic metric sector. That earlier symbol is
better called $\mathscr D_{\mathcal Q}$ here. Its effect is
already carried by the $-\lambda(D_s\mathcal Q)^*$ term in
\emph{each} of the three channels above. The old and new conventions
must not be added together or counted twice.

In the fixed-coframe, $\lambda=0$ specialization and $c=1$,
write $\mathcal R^{ab}:=T^a{}_{\mu}T^b{}_{\nu}
(G^{\mu\nu}+\Lambda g^{\mu\nu})$ and use similarly
$T_{\mathrm{cl/q}}^{ab}$ for the frame-projected stress. Then
\begin{equation}\label{eq:qgd141-explicit}
\boxed{\begin{aligned}
\mathscr G_{(s)}^{\,a}&=-\frac{s_s}{8\pi G_N}
\mathcal R^{ab}\sigma_b^{(s)},\\
\mathscr T_{(s)}^{\,a}&=s_s T_{\mathrm{cl}}^{ab}\sigma_b^{(s)},\\
\mathscr Q_{(s)}^{\,a}&=s_s T_{\mathrm{q}}^{ab}\sigma_b^{(s)},\\
\mathscr S_{(s)}^{\,a}&=\mathscr G_{(s)}^{\,a}
+\mathscr T_{(s)}^{\,a}+\mathscr Q_{(s)}^{\,a}=0.
\end{aligned}}
\end{equation}
\paragraph{Compact natural-invariant projection of the Einstein residual.}
Equivalently, adding the three channels in
Eq.~\eqref{eq:qgd141-explicit} gives the concise field equation
\begin{equation}\label{eq:qgd141-projected-compact}
\boxed{\displaystyle
s_s\sigma_b^{(s)}\left[
-\frac{\mathcal R^{ab}}{8\pi G_N}
+T_{\mathrm{cl}}^{ab}+T_{\mathrm q}^{ab}\right]=0,
\qquad(s=1,\ldots,N).}
\end{equation}
This is the exact action-derived \emph{projected Einstein residual}
in the stated fixed-coframe, $\lambda\mathcal Q=0$, metric-only
specialization, in units $c=1$.  The curvature symbol
$\mathcal R^{ab}$ is the frame pullback of the Einstein tensor
\emph{plus} the cosmological term, not the Ricci tensor alone.
Under a regular surjective metric map, the full collection of sector
equations is equivalent to the semiclassical Einstein equation;
on singular or rank-deficient sector branches, vanishing projections
need not imply vanishing of the full Einstein residual.
For a nonzero specified derivative slot $\lambda\mathcal Q_{ab}$,
Eq.~\eqref{eq:qgd141-projected-compact} is \emph{not} the complete
Euler equation: the volume-weighted functional-adjoint contribution
in Eq.~\eqref{eq:qgd137-sector} must also be retained.
Restoring units replaces $8\pi G_N$ with $8\pi G_N/c^4$ in
the Einstein source equation. For the full derivative map, the
$-\lambda(D_s\mathcal Q)^*$ term acts on \emph{each} $\Pi_X$;
it is not an additional independently propagating quantum source.
This reproduces Option C, because
$2s_s\Pi^{ab}\sigma_b^{(s)}=
\lambda[(D_s\mathcal Q)^*_{\sqrt{-g}}\Pi]^a$ is exactly
Eq.~\eqref{eq:qgd141-balance} rearranged.


\paragraph{Full-master-metric extension: the projected Einstein residual with derivative response.}
The compact projection in Eq.~\eqref{eq:qgd141-projected-compact}
is valid only for $\lambda\mathcal Q=0$. For a specified differentiable
$\mathcal Q_{ab}[\sigma;T]$, the \emph{complete} stationarity condition
is instead obtained by retaining its functional-adjoint term. Define
$\Pi^{ab}:=\Pi_{\mathsf G}^{ab}+\Pi_{\mathsf T}^{ab}+\Pi_{\mathsf Q}^{ab}$,
where the metric Euler tensors $\Pi_X^{ab}$ are defined in
Eq.~\eqref{eq:qgd141-channels}; in $c=1$ units this is
$\Pi^{ab}=\frac12[-\mathcal R^{ab}/(8\pi G_N)
+T_{\mathrm{cl}}^{ab}+T_{\mathrm q}^{ab}]$.
The Option-C form of the complete natural-invariant equation is therefore
\begin{equation}\label{eq:qgd141-full-projected-einstein}
\boxed{\begin{aligned}
&s_s\sigma_b^{(s)}\left[
-\frac{\mathcal R^{ab}}{8\pi G_N}
+T_{\mathrm{cl}}^{ab}+T_{\mathrm q}^{ab}
\right]\\[4pt]
&\qquad=\lambda\left[
(D_s\mathcal Q)^*_{\sqrt{-g}}\Pi\right]^a.
\end{aligned}}
\end{equation}
Here $s_s=-\varepsilon_s$, $D_s\mathcal Q$ is the Fr\'echet
variation with respect to $\sigma^{(s)}$ at fixed coframe, and
$\mathcal R^{ab}:=T^a{}_{\mu}T^b{}_{\nu}
(G^{\mu\nu}+\Lambda g^{\mu\nu})$ is \emph{not} the Ricci tensor.
Both sides arise by varying the same \emph{metric-only} renormalized
action through the full composite metric; the right side does not
constitute an independently postulated source. The matter-loop
contribution $T_{\mathrm q}^{ab}$ and the historical metric operator
$\mathcal Q_{ab}$ remain distinct. A nonzero $\lambda$ is not
microscopically predicted by this rewriting, and physical equivalence
to the metric field equations requires an admissible full-rank metric map.
Setting $\lambda=0$ returns Eq.~\eqref{eq:qgd141-projected-compact}
without an extra assumption.

\begin{proposition}[Regular full-rank correspondence and its limit]
\label{prop:qgd141-rank}
For a metric-only action on a regular branch where the full
$D_\sigma g$ is surjective on admissible metric variations,
$\mathscr S_{(s)}^{\,a}=0$ for all sectors is equivalent to
$\mathcal E_g^{\mu\nu}=0$. The sum of the three channels is then zero;
the channels individually need not vanish. There is no independent
$\sigma$-wave energy or new force coefficient obtained just by
splitting the zero residual.
\end{proposition}
\begin{proof}
Use the complete functional chain rule
$\delta\Gamma=\int\sqrt{-g}\,\mathcal E_g^{\mu\nu}
(D_\sigma g)_{\mu\nu}\delta\sigma$.
Surjectivity gives the converse; linearity gives the channel identity.
\end{proof}

\subsection{The balance is not an independent local energy-conservation law}
The geometric Bianchi identity and renormalized matter Ward identity
imply, on the matter equations and with appropriate state-dependent
backreaction, $\nabla_\mu(T_{\mathrm{cl}}^{\mu\nu}
+T_{\mathrm{q}}^{\mu\nu})=0$; they do not imply separate
conservation of every channel. Microscopic reference-space Noether
conservation (Chapter~6), the projected Euler equation, and
ADM/Bondi/quasilocal energies are \emph{three different constructions}.
This distinction prevents reading $\mathscr G,\mathscr T,\mathscr Q$
as three gauge-invariant local energy densities.

\textbf{Renormalization caveat.} A finite local counterterm may be
moved between $\Gamma_{\mathsf G}$ and $\Gamma_{\mathsf Q}$,
with an equal and opposite change, leaving $\Gamma_{\mathrm{metric}}$
and $\mathscr S$ unchanged. The individual named channels therefore
require a declared subtraction convention; only the total field
equation is scheme independent after physical parameters are matched.

\subsection{A nontrivial test of both quantum energy and pressure}
For the published free massless-Dirac Boulware far field of
Section~\ref{sec:quantumstructure}, use $G_N=c=\hbar=1$ and
$\langle T^\mu{}_{\nu}\rangle=M\operatorname{diag}(4,2,-3,-3)/(40\pi^2r^5)$.
The effective energy density and radial pressure are respectively
$\epsilon_q=-M/(10\pi^2r^5)$ and
$p_{r,q}=M/(20\pi^2r^5)$.
The spherical Einsteinian source equations give, at leading quantum
far-field order,
\begin{equation}\label{eq:qgd141-radial}
\boxed{\begin{aligned}
\delta m'(r)&=4\pi r^2\epsilon_q=-\frac{2M}{5\pi r^3},
&\delta m(r)&=\frac{M}{5\pi r^2},\\
\frac{\dd}{\dd r}\ln D
&=\frac{8\pi r(\epsilon_q+p_{r,q})}{1-2m/r}
=-\frac{2M}{5\pi r^4}+\text{higher orders},
&\ln D&=\frac{2M}{15\pi r^3}+\text{higher orders},
\end{aligned}}
\end{equation}
where $D=1+B^2-A^2$. The first equation measures the
\emph{enclosed energy} response; the second responds to
\emph{energy density plus radial pressure}, and explains why
$A^2\ne B^2$ for this quantum state. Both arise from the same
$T_q^{\mu\nu}$ and the same Euler balance, rather than independent
historical $Q_\mu$ and $G_\mu$ source guesses.

\textbf{Conclusion.} The historical formula retains conceptual
value as an architecture of gravity, matter and quantum response.
Its justified successor is Eq.~\eqref{eq:qgd141-balance},
not $\Box_g\sigma_\mu=Q_\mu+G_\mu+T_\mu$ unless a new,
properly normalized and stable independent covector kinetic action
is separately established. This is a completed action-based
reinterpretation, not a microscopic proof of the historical wave law.
% ===== END QGD-141 =====


\clearpage
\end{document}
